% Options for packages loaded elsewhere
\PassOptionsToPackage{unicode}{hyperref}
\PassOptionsToPackage{hyphens}{url}
\documentclass[
]{article}
\usepackage{xcolor}
\usepackage{amsmath,amssymb}
\setcounter{secnumdepth}{-\maxdimen} % remove section numbering
\usepackage{iftex}
\ifPDFTeX
  \usepackage[T1]{fontenc}
  \usepackage[utf8]{inputenc}
  \usepackage{textcomp} % provide euro and other symbols
\else % if luatex or xetex
  \usepackage{unicode-math} % this also loads fontspec
  \defaultfontfeatures{Scale=MatchLowercase}
  \defaultfontfeatures[\rmfamily]{Ligatures=TeX,Scale=1}
\fi
\usepackage{lmodern}
\ifPDFTeX\else
  % xetex/luatex font selection
\fi
% Use upquote if available, for straight quotes in verbatim environments
\IfFileExists{upquote.sty}{\usepackage{upquote}}{}
\IfFileExists{microtype.sty}{% use microtype if available
  \usepackage[]{microtype}
  \UseMicrotypeSet[protrusion]{basicmath} % disable protrusion for tt fonts
}{}
\makeatletter
\@ifundefined{KOMAClassName}{% if non-KOMA class
  \IfFileExists{parskip.sty}{%
    \usepackage{parskip}
  }{% else
    \setlength{\parindent}{0pt}
    \setlength{\parskip}{6pt plus 2pt minus 1pt}}
}{% if KOMA class
  \KOMAoptions{parskip=half}}
\makeatother
\usepackage{color}
\usepackage{fancyvrb}
\newcommand{\VerbBar}{|}
\newcommand{\VERB}{\Verb[commandchars=\\\{\}]}
\DefineVerbatimEnvironment{Highlighting}{Verbatim}{commandchars=\\\{\}}
% Add ',fontsize=\small' for more characters per line
\newenvironment{Shaded}{}{}
\newcommand{\AlertTok}[1]{\textcolor[rgb]{1.00,0.00,0.00}{\textbf{#1}}}
\newcommand{\AnnotationTok}[1]{\textcolor[rgb]{0.38,0.63,0.69}{\textbf{\textit{#1}}}}
\newcommand{\AttributeTok}[1]{\textcolor[rgb]{0.49,0.56,0.16}{#1}}
\newcommand{\BaseNTok}[1]{\textcolor[rgb]{0.25,0.63,0.44}{#1}}
\newcommand{\BuiltInTok}[1]{\textcolor[rgb]{0.00,0.50,0.00}{#1}}
\newcommand{\CharTok}[1]{\textcolor[rgb]{0.25,0.44,0.63}{#1}}
\newcommand{\CommentTok}[1]{\textcolor[rgb]{0.38,0.63,0.69}{\textit{#1}}}
\newcommand{\CommentVarTok}[1]{\textcolor[rgb]{0.38,0.63,0.69}{\textbf{\textit{#1}}}}
\newcommand{\ConstantTok}[1]{\textcolor[rgb]{0.53,0.00,0.00}{#1}}
\newcommand{\ControlFlowTok}[1]{\textcolor[rgb]{0.00,0.44,0.13}{\textbf{#1}}}
\newcommand{\DataTypeTok}[1]{\textcolor[rgb]{0.56,0.13,0.00}{#1}}
\newcommand{\DecValTok}[1]{\textcolor[rgb]{0.25,0.63,0.44}{#1}}
\newcommand{\DocumentationTok}[1]{\textcolor[rgb]{0.73,0.13,0.13}{\textit{#1}}}
\newcommand{\ErrorTok}[1]{\textcolor[rgb]{1.00,0.00,0.00}{\textbf{#1}}}
\newcommand{\ExtensionTok}[1]{#1}
\newcommand{\FloatTok}[1]{\textcolor[rgb]{0.25,0.63,0.44}{#1}}
\newcommand{\FunctionTok}[1]{\textcolor[rgb]{0.02,0.16,0.49}{#1}}
\newcommand{\ImportTok}[1]{\textcolor[rgb]{0.00,0.50,0.00}{\textbf{#1}}}
\newcommand{\InformationTok}[1]{\textcolor[rgb]{0.38,0.63,0.69}{\textbf{\textit{#1}}}}
\newcommand{\KeywordTok}[1]{\textcolor[rgb]{0.00,0.44,0.13}{\textbf{#1}}}
\newcommand{\NormalTok}[1]{#1}
\newcommand{\OperatorTok}[1]{\textcolor[rgb]{0.40,0.40,0.40}{#1}}
\newcommand{\OtherTok}[1]{\textcolor[rgb]{0.00,0.44,0.13}{#1}}
\newcommand{\PreprocessorTok}[1]{\textcolor[rgb]{0.74,0.48,0.00}{#1}}
\newcommand{\RegionMarkerTok}[1]{#1}
\newcommand{\SpecialCharTok}[1]{\textcolor[rgb]{0.25,0.44,0.63}{#1}}
\newcommand{\SpecialStringTok}[1]{\textcolor[rgb]{0.73,0.40,0.53}{#1}}
\newcommand{\StringTok}[1]{\textcolor[rgb]{0.25,0.44,0.63}{#1}}
\newcommand{\VariableTok}[1]{\textcolor[rgb]{0.10,0.09,0.49}{#1}}
\newcommand{\VerbatimStringTok}[1]{\textcolor[rgb]{0.25,0.44,0.63}{#1}}
\newcommand{\WarningTok}[1]{\textcolor[rgb]{0.38,0.63,0.69}{\textbf{\textit{#1}}}}
\usepackage{longtable,booktabs,array}
\newcounter{none} % for unnumbered tables
\usepackage{calc} % for calculating minipage widths
% Correct order of tables after \paragraph or \subparagraph
\usepackage{etoolbox}
\makeatletter
\patchcmd\longtable{\par}{\if@noskipsec\mbox{}\fi\par}{}{}
\makeatother
% Allow footnotes in longtable head/foot
\IfFileExists{footnotehyper.sty}{\usepackage{footnotehyper}}{\usepackage{footnote}}
\makesavenoteenv{longtable}
\setlength{\emergencystretch}{3em} % prevent overfull lines
\providecommand{\tightlist}{%
  \setlength{\itemsep}{0pt}\setlength{\parskip}{0pt}}
\usepackage{bookmark}
\IfFileExists{xurl.sty}{\usepackage{xurl}}{} % add URL line breaks if available
\urlstyle{same}
\hypersetup{
  hidelinks,
  pdfcreator={LaTeX via pandoc}}

\author{}
\date{}

\begin{document}

\textit{[Image omitted: images/d3d8b821daec3e2e0c83d6c64b6477d2fcdfb88e9fd669608a0d57dcba07406f.jpg]}

\section{usenix®}\label{usenix}

THE ADVANCED

COMPUTING SYSTEMS

ASSOCIATION

\section{Ray: A Distributed Framework for Emerging AI
Applications}\label{ray-a-distributed-framework-for-emerging-ai-applications}

Philipp Moritz, Robert Nishihara, Stephanie Wang, Alexey Tumanov,
Richard Liaw, Eric Liang, Melih Elibol, Zongheng Yang, William Paul,
Michael I. Jordan, and Ion Stoica, UC Berkeley

https://www.usenix.org/conference/osdi18/presentation/nishihara

This paper is included in the Proceedings of the 13th USENIX Symposium
on Operating Systems Design and Implementation (OSDI '18).

October 8--10, 2018 • Carlsbad, CA, USA

ISBN 978-1-939133-08-3

Open access to the Proceedings of the 13th USENIX Symposium on Operating
Systems Design and Implementation is sponsored by USENIX.

\section{Ray: A Distributed Framework for Emerging AI
Applications}\label{ray-a-distributed-framework-for-emerging-ai-applications-1}

Philipp Moritz*, Robert Nishihara*, Stephanie Wang, Alexey Tumanov,
Richard Liaw, Eric Liang, Melih Elibol, Zongheng Yang, William Paul,
Michael I. Jordan, Ion Stoica University of California, Berkeley

\section{Abstract}\label{abstract}

The next generation of AI applications will continuously interact with
the environment and learn from these interactions. These applications
impose new and demanding systems requirements, both in terms of
performance and flexibility. In this paper, we consider these
requirements and present Ray---a distributed system to address them. Ray
implements a unified interface that can express both task-parallel and
actor-based computations, supported by a single dynamic execution
engine. To meet the performance requirements, Ray employs a distributed
scheduler and a distributed and fault-tolerant store to manage the
system's control state. In our experiments, we demonstrate scaling
beyond 1.8 million tasks per second and better performance than existing
specialized systems for several challenging reinforcement learning
applications.

\section{1 Introduction}\label{introduction}

Over the past two decades, many organizations have been collecting---and
aiming to exploit---ever-growing quantities of data. This has led to the
development of a plethora of frameworks for distributed data analysis,
including batch {[}20, 64, 28{]}, streaming {[}15, 39, 31{]}, and graph
{[}34, 35, 24{]} processing systems. The success of these frameworks has
made it possible for organizations to analyze large data sets as a core
part of their business or scientific strategy, and has ushered in the
age of ``Big Data.''

More recently, the scope of data-focused applications has expanded to
encompass more complex artificial intelligence (AI) or machine learning
(ML) techniques {[}30{]}. The paradigm case is that of supervised
learning, where data points are accompanied by labels, and where the
workhorse technology for mapping data points to labels is provided by
deep neural networks. The complexity of these deep networks has led to
another flurry of frameworks that focus on the training of deep neural
networks and their use in prediction. These frameworks often leverage
specialized hardware (e.g., GPUs and TPUs), with the goal of reducing
training time in a batch setting. Examples include TensorFlow {[}7{]},
MXNet {[}18{]}, and PyTorch {[}46{]}.

The promise of AI is, however, far broader than classical supervised
learning. Emerging AI applications must increasingly operate in dynamic
environments, react to changes in the environment, and take sequences of
actions to accomplish long-term goals {[}8, 43{]}. They must aim not
only to exploit the data gathered, but also to explore the space of
possible actions. These broader requirements are naturally framed within
the paradigm of reinforcement learning (RL). RL deals with learning to
operate continuously within an uncertain environment based on delayed
and limited feedback {[}56{]}. RL-based systems have already yielded
remarkable results, such as Google's AlphaGo beating a human world
champion {[}54{]}, and are beginning to find their way into dialogue
systems, UAVs {[}42{]}, and robotic manipulation {[}25, 60{]}.

The central goal of an RL application is to learn a policy---a mapping
from the state of the environment to a choice of action---that yields
effective performance over time, e.g., winning a game or piloting a
drone. Finding effective policies in large-scale applications requires
three main capabilities. First, RL methods often rely on simulation to
evaluate policies. Simulations make it possible to explore many
different choices of action sequences and to learn about the long-term
consequences of those choices. Second, like their supervised learning
counterparts, RL algorithms need to perform distributed training to
improve the policy based on data generated through simulations or
interactions with the physical environment. Third, policies are intended
to provide solutions to control problems, and thus it is necessary to
serve the policy in interactive closed-loop and open-loop control
scenarios.

These characteristics drive new systems requirements: a system for RL
must support fine-grained computations (e.g., rendering actions in
milliseconds when interacting with the real world, and performing vast
numbers of sim-

ulations), must support heterogeneity both in time (e.g., a simulation
may take milliseconds or hours) and in resource usage (e.g., GPUs for
training and CPUs for simulations), and must support dynamic execution,
as results of simulations or interactions with the environment can
change future computations. Thus, we need a dynamic computation
framework that handles millions of heterogeneous tasks per second at
millisecond-level latencies.

Existing frameworks that have been developed for Big Data workloads or
for supervised learning workloads fall short of satisfying these new
requirements for RL. Bulk-synchronous parallel systems such as MapReduce
{[}20{]}, Apache Spark {[}64{]}, and Dryad {[}28{]} do not support
fine-grained simulation or policy serving. Task-parallel systems such as
CIEL {[}40{]} and Dask {[}48{]} provide little support for distributed
training and serving. The same is true for streaming systems such as
Naiad {[}39{]} and Storm {[}31{]}. Distributed deep-learning frameworks
such as TensorFlow {[}7{]} and MXNet {[}18{]} do not naturally support
simulation and serving. Finally, model-serving systems such as
TensorFlow Serving {[}6{]} and Clipper {[}19{]} support neither training
nor simulation.

While in principle one could develop an end-to-end solution by stitching
together several existing systems (e.g., Horovod {[}53{]} for
distributed training, Clipper {[}19{]} for serving, and CIEL {[}40{]}
for simulation), in practice this approach is untenable due to the tight
coupling of these components within applications. As a result,
researchers and practitioners today build one-off systems for
specialized RL applications {[}58, 41, 54, 44, 49, 5{]}. This approach
imposes a massive systems engineering burden on the development of
distributed applications by essentially pushing standard systems
challenges like scheduling, fault tolerance, and data movement onto each
application.

In this paper, we propose Ray, a general-purpose cluster-computing
framework that enables simulation, training, and serving for RL
applications. The requirements of these workloads range from lightweight
and stateless computations, such as for simulation, to long-running and
stateful computations, such as for training. To satisfy these
requirements, Ray implements a unified interface that can express both
task-parallel and actor-based computations. Tasks enable Ray to
efficiently and dynamically load balance simulations, process large
inputs and state spaces (e.g., images, video), and recover from
failures. In contrast, actors enable Ray to efficiently support stateful
computations, such as model training, and expose shared mutable state to
clients, (e.g., a parameter server). Ray implements the actor and the
task abstractions on top of a single dynamic execution engine that is
highly scalable and fault tolerant.

To meet the performance requirements, Ray distributes two components
that are typically centralized in existing frameworks {[}64, 28, 40{]}:
(1) the task scheduler and (2) a metadata store which maintains the
computation lineage and a directory for data objects. This allows Ray to
schedule millions of tasks per second with millisecond-level latencies.
Furthermore, Ray provides lineage-based fault tolerance for tasks and
actors, and replication-based fault tolerance for the metadata store.

\textit{[Image omitted: images/9bcab303d552167943db871c3b8febcf7320127759d9cf39adc9e4110b5070c3.jpg]}

flowchart

\begin{Shaded}
\begin{Highlighting}[]
\NormalTok{graph TD}
\NormalTok{    A["Training Policy improvement (e.g., SGD)"] {-}{-}\textgreater{}|policy| B["Serving Policy evaluation"]}
\NormalTok{    B {-}{-}\textgreater{}|action (a\_i)| C["Simulation"]}
\NormalTok{    C {-}{-}\textgreater{}|state (s\_\{i+1\}) (observation)| B}
\NormalTok{    B {-}{-}\textgreater{}|reward (r\_\{i+1\})| C}
\NormalTok{    C {-}{-}\textgreater{}|trajectory: s\_0, (s\_1, r\_1), ..., (s\_n, r\_n)| A}
\NormalTok{    style A fill:\#f9f,stroke:\#333}
\NormalTok{    style B fill:\#ccf,stroke:\#333}
\NormalTok{    style C fill:\#cfc,stroke:\#333}
\end{Highlighting}
\end{Shaded}

Figure 1: Example of an RL system.

While Ray supports serving, training, and simulation in the context of
RL applications, this does not mean that it should be viewed as a
replacement for systems that provide solutions for these workloads in
other contexts. In particular, Ray does not aim to substitute for
serving systems like Clipper {[}19{]} and TensorFlow Serving {[}6{]}, as
these systems address a broader set of challenges in deploying models,
including model management, testing, and model composition. Similarly,
despite its flexibility, Ray is not a substitute for generic
data-parallel frameworks, such as Spark {[}64{]}, as it currently lacks
the rich functionality and APIs (e.g., straggler mitigation, query
optimization) that these frameworks provide.

We make the following contributions:

\begin{itemize}
\tightlist
\item
  We design and build the first distributed framework that unifies
  training, simulation, and serving---necessary components of emerging
  RL applications.\\
\item
  To support these workloads, we unify the actor and task-parallel
  abstractions on top of a dynamic task execution engine.\\
\item
  To achieve scalability and fault tolerance, we propose a system design
  principle in which control state is stored in a sharded metadata store
  and all other system components are stateless.\\
\item
  To achieve scalability, we propose a bottom-up distributed scheduling
  strategy.
\end{itemize}

\section{2 Motivation and
Requirements}\label{motivation-and-requirements}

We begin by considering the basic components of an RL system and
fleshing out the key requirements for Ray. As shown in Figure 1, in an
RL setting, an agent interacts repeatedly with the environment. The goal
of the agent is to learn a policy that maximizes a reward. A policy is

\begin{Shaded}
\begin{Highlighting}[]
\OperatorTok{//}\NormalTok{ evaluate policy by interacting }\ControlFlowTok{with}\NormalTok{ env. (e.g., simulator)}
\NormalTok{rollout(policy, environment):}
\NormalTok{trajectory }\OperatorTok{=}\NormalTok{ []}
\NormalTok{state }\OperatorTok{=}\NormalTok{ environment.initial\_state()}
\ControlFlowTok{while}\NormalTok{ (}\KeywordTok{not}\NormalTok{ environment.has\_terminated()):}
\NormalTok{    action }\OperatorTok{=}\NormalTok{ policy.compute(state) }\OperatorTok{//}\NormalTok{ Serving}
\NormalTok{    state, reward }\OperatorTok{=}\NormalTok{ environment.step(action) }\OperatorTok{//}\NormalTok{ Simulation}
\NormalTok{    trajectory.append(state, reward)}
\ControlFlowTok{return}\NormalTok{ trajectory }
\end{Highlighting}
\end{Shaded}

\begin{Shaded}
\begin{Highlighting}[]
\OperatorTok{//}\NormalTok{ improve policy iteratively until it converges}
\NormalTok{train\_policy(environment):}
\NormalTok{    policy }\OperatorTok{=}\NormalTok{ initial\_policy()}
    \ControlFlowTok{while}\NormalTok{ (policy has }\KeywordTok{not}\NormalTok{ converged):}
\NormalTok{    trajectories }\OperatorTok{=}\NormalTok{ []}
    \ControlFlowTok{for}\NormalTok{ i }\ImportTok{from} \DecValTok{1}\NormalTok{ to k:}
    \OperatorTok{//}\NormalTok{ evaluate policy by generating k rollouts}
\NormalTok{    trajectories.append(rollout(policy, environment))}
    \OperatorTok{//}\NormalTok{ improve policy}
\NormalTok{    policy }\OperatorTok{=}\NormalTok{ policy.update(trajectories) }\OperatorTok{//}\NormalTok{ Training}
\ControlFlowTok{return}\NormalTok{ policy }
\end{Highlighting}
\end{Shaded}

Figure 2: Typical RL pseudocode for learning a policy.

a mapping from the state of the environment to a choice of action. The
precise definitions of environment, agent, state, action, and reward are
application-specific.

To learn a policy, an agent typically employs a two-step process: (1)
policy evaluation and (2) policy improvement. To evaluate the policy,
the agent interacts with the environment (e.g., with a simulation of the
environment) to generate trajectories, where a trajectory consists of a
sequence of (state, reward) tuples produced by the current policy. Then,
the agent uses these trajectories to improve the policy; i.e., to update
the policy in the direction of the gradient that maximizes the reward.
Figure 2 shows an example of the pseudocode used by an agent to learn a
policy. This pseudocode evaluates the policy by invoking
rollout(environment, policy) to generate trajectories. train\_policy()
then uses these trajectories to improve the current policy via
policy.update(trajectories). This process repeats until the policy
converges.

Thus, a framework for RL applications must provide efficient support for
training, serving, and simulation (Figure 1). Next, we briefly describe
these workloads.

Training typically involves running stochastic gradient descent (SGD),
often in a distributed setting, to update the policy. Distributed SGD
typically relies on an allreduce aggregation step or a parameter server
{[}32{]}.

Serving uses the trained policy to render an action based on the current
state of the environment. A serving system aims to minimize latency, and
maximize the number of decisions per second. To scale, load is typically
balanced across multiple nodes serving the policy.

Finally, most existing RL applications use simulations to evaluate the
policy---current RL algorithms are not sample-efficient enough to rely
solely on data obtained from interactions with the physical world. These
simulations vary widely in complexity. They might take a few ms (e.g.,
simulate a move in a chess game) to minutes (e.g., simulate a realistic
environment for a self-driving car).

In contrast with supervised learning, in which training and serving can
be handled separately by different systems, in RL all three of these
workloads are tightly coupled in a single application, with stringent
latency requirements between them. Currently, no framework supports this
coupling of workloads. In theory, multiple specialized frameworks could
be stitched together to provide the overall capabilities, but in
practice, the resulting data movement and latency between systems is
prohibitive in the context of RL. As a result, researchers and
practitioners have been building their own one-off systems.

This state of affairs calls for the development of new distributed
frameworks for RL that can efficiently support training, serving, and
simulation. In particular, such a framework should satisfy the following
requirements:

Fine-grained, heterogeneous computations. The duration of a computation
can range from milliseconds (e.g., taking an action) to hours (e.g.,
training a complex policy). Additionally, training often requires
heterogeneous hardware (e.g., CPUs, GPUs, or TPUs).

Flexible computation model. RL applications require both stateless and
stateful computations. Stateless computations can be executed on any
node in the system, which makes it easy to achieve load balancing and
movement of computation to data, if needed. Thus stateless computations
are a good fit for fine-grained simulation and data processing, such as
extracting features from images or videos. In contrast stateful
computations are a good fit for implementing parameter servers,
performing repeated computation on GPU-backed data, or running
third-party simulators that do not expose their state.

Dynamic execution. Several components of RL applications require dynamic
execution, as the order in which computations finish is not always known
in advance (e.g., the order in which simulations finish), and the
results of a computation can determine future computations (e.g., the
results of a simulation will determine whether we need to perform more
simulations).

We make two final comments. First, to achieve high utilization in large
clusters, such a framework must handle millions of tasks per second.*
Second, such a framework is not intended for implementing deep neural
networks or complex simulators from scratch. Instead, it should enable
seamless integration with existing simulators {[}13, 11, 59{]} and deep
learning frameworks {[}7, 18, 46, 29{]}.

Name

Description

futures = f.remote(args)

Execute function f remotely. f.remote() can take objects or futures as
inputs and returns one or more futures. This is non-blocking.

objects = ray.get(futures)

Return the values associated with one or more futures. This is blocking.

ready\_futures = ray.wait(futures, k, timeout)

Return the futures whose corresponding tasks have completed as soon as
either k have completed or the timeout expires.

actor = Class.remote(args)futures = actor.method.remote(args)

Instantiate class Class as a remote actor, and return a handle to it.
Call a method on the remote actor and return one or more futures. Both
are non-blocking.

Table 1: Ray API

\section{3 Programming and Computation
Model}\label{programming-and-computation-model}

Ray implements a dynamic task graph computation model, i.e., it models
an application as a graph of dependent tasks that evolves during
execution. On top of this model, Ray provides both an actor and a
task-parallel programming abstraction. This unification differentiates
Ray from related systems like CIEL, which only provides a task-parallel
abstraction, and from Orleans {[}14{]} or Akka {[}1{]}, which primarily
provide an actor abstraction.

\section{3.1 Programming Model}\label{programming-model}

Tasks. A task represents the execution of a remote function on a
stateless worker. When a remote function is invoked, a future
representing the result of the task is returned immediately. Futures can
be retrieved using ray.get() and passed as arguments into other remote
functions without waiting for their result. This allows the user to
express parallelism while capturing data dependencies. Table 1 shows
Ray's API.

Remote functions operate on immutable objects and are expected to be
stateless and side-effect free: their outputs are determined solely by
their inputs. This implies idempotence, which simplifies fault tolerance
through function re-execution on failure.

Actors. An actor represents a stateful computation. Each actor exposes
methods that can be invoked remotely and are executed serially. A method
execution is similar to a task, in that it executes remotely and returns
a future, but differs in that it executes on a stateful worker. A handle
to an actor can be passed to other actors or tasks, making it possible
for them to invoke methods on that actor.

Tasks (stateless)

Actors (stateful)

Fine-grained load balancing

Coarse-grained load balancing

Support for object locality

Poor locality support

High overhead for small updates

Low overhead for small updates

Efficient failure handling

Overhead from checkpointing

Table 2: Tasks vs.~actors tradeoffs.

Table 2 summarizes the properties of tasks and actors. Tasks enable
fine-grained load balancing through leveraging load-aware scheduling at
task granularity, input data locality, as each task can be scheduled on
the node storing its inputs, and low recovery overhead, as there is no
need to checkpoint and recover intermediate state. In contrast, actors
provide much more efficient fine-grained updates, as these updates are
performed on internal rather than external state, which typically
requires serialization and deserialization. For example, actors can be
used to implement parameter servers {[}32{]} and GPU-based iterative
computations (e.g., training). In addition, actors can be used to wrap
third-party simulators and other opaque handles that are hard to
serialize.

To satisfy the requirements for heterogeneity and flexibility (Section
2), we augment the API in three ways. First, to handle concurrent tasks
with heterogeneous durations, we introduce ray.wait(), which waits for
the first k available results, instead of waiting for all results like
ray.get(). Second, to handle resource-heterogeneous tasks, we enable
developers to specify resource requirements so that the Ray scheduler
can efficiently manage resources. Third, to improve flexibility, we
enable nested remote functions, meaning that remote functions can invoke
other remote functions. This is also critical for achieving high
scalability (Section 4), as it enables multiple processes to invoke
remote functions in a distributed fashion.

\section{3.2 Computation Model}\label{computation-model}

Ray employs a dynamic task graph computation model {[}21{]}, in which
the execution of both remote functions and actor methods is
automatically triggered by the system when their inputs become
available. In this section, we describe how the computation graph
(Figure 4) is constructed from a user program (Figure 3). This program
uses the API in Table 1 to implement the pseudocode from Figure 2.

Ignoring actors first, there are two types of nodes in a computation
graph: data objects and remote function invocations, or tasks. There are
also two types of edges: data edges and control edges. Data edges
capture the de-

\begin{Shaded}
\begin{Highlighting}[]
\AttributeTok{@ray.remote}
\KeywordTok{def}\NormalTok{ create\_policy():}
    \CommentTok{\# Initialize the policy randomly.}
    \ControlFlowTok{return}\NormalTok{ policy}

\AttributeTok{@ray.remote}\NormalTok{(num\_gpus}\OperatorTok{=}\DecValTok{1}\NormalTok{)}
\KeywordTok{class}\NormalTok{ Simulator(}\BuiltInTok{object}\NormalTok{):}
    \KeywordTok{def} \FunctionTok{\_\_init\_\_}\NormalTok{(}\VariableTok{self}\NormalTok{):}
    \CommentTok{\# Initialize the environment.}
    \VariableTok{self}\NormalTok{.env }\OperatorTok{=}\NormalTok{ Environment()}
    \KeywordTok{def}\NormalTok{ rollout(}\VariableTok{self}\NormalTok{, policy, num\_steps):}
\NormalTok{    observations }\OperatorTok{=}\NormalTok{ []}
\NormalTok{    observation }\OperatorTok{=} \VariableTok{self}\NormalTok{.env.current\_state()}
    \ControlFlowTok{for}\NormalTok{ \_ }\KeywordTok{in} \BuiltInTok{range}\NormalTok{(num\_steps):}
\NormalTok{    action }\OperatorTok{=}\NormalTok{ policy(observation)}
\NormalTok{    observation }\OperatorTok{=} \VariableTok{self}\NormalTok{.env.step(action)}
\NormalTok{    observations.append(observation)}
    \ControlFlowTok{return}\NormalTok{ observations}

\AttributeTok{@ray.remote}\NormalTok{(num\_gpus}\OperatorTok{=}\DecValTok{2}\NormalTok{)}
\KeywordTok{def}\NormalTok{ update\_policy(policy, }\OperatorTok{*}\NormalTok{rollouts):}
    \CommentTok{\# Update the policy.}
    \ControlFlowTok{return}\NormalTok{ policy}

\AttributeTok{@ray.remote}
\KeywordTok{def}\NormalTok{ train\_policy():}
    \CommentTok{\# Create a policy.}
\NormalTok{    policy\_id }\OperatorTok{=}\NormalTok{ create\_policy.remote()}
    \CommentTok{\# Create 10 actors.}
\NormalTok{    simulators }\OperatorTok{=}\NormalTok{ [Simulator.remote() }\ControlFlowTok{for}\NormalTok{ \_ }\KeywordTok{in} \BuiltInTok{range}\NormalTok{(}\DecValTok{10}\NormalTok{)]}
    \CommentTok{\# Do 100 steps of training.}
    \ControlFlowTok{for}\NormalTok{ \_ }\KeywordTok{in} \BuiltInTok{range}\NormalTok{(}\DecValTok{100}\NormalTok{):}
    \CommentTok{\# Perform one rollout on each actor.}
\NormalTok{    rollout\_ids }\OperatorTok{=}\NormalTok{ [s. rollout.remote(policy\_id)}
    \ControlFlowTok{for}\NormalTok{ s }\KeywordTok{in}\NormalTok{ simulators]}
    \CommentTok{\# Update the policy with the rollouts.}
\NormalTok{    policy\_id }\OperatorTok{=}\NormalTok{ update\_policy.remote(policy\_id, }\OperatorTok{*}\NormalTok{rollouts\_ids)}
    \ControlFlowTok{return}\NormalTok{ ray.get(policy\_id) }
\end{Highlighting}
\end{Shaded}

Figure 3: Python code implementing the example in Figure 2 in Ray. Note
that @ray.remote indicates remote functions and actors. Invocations of
remote functions and actor methods return futures, which can be passed
to subsequent remote functions or actor methods to encode task
dependencies. Each actor has an environment object self.env shared
between all of its methods.

pendencies between data objects and tasks. More precisely, if data
object D is an output of task T, we add a data edge from T to D.
Similarly, if D is an input to T, we add a data edge from D to T.
Control edges capture the computation dependencies that result from
nested remote functions (Section 3.1): if task \(T_{1}\) invokes task
\(T_{2}\) , then we add a control edge from \(T_{1}\) to \(T_{2}\) .

Actor method invocations are also represented as nodes in the
computation graph. They are identical to tasks with one key difference.
To capture the state dependency across subsequent method invocations on
the same actor, we add a third type of edge: a stateful edge. If method
\(M_{j}\) is called right after method \(M_{i}\) on the same actor, then
we add a stateful edge from \(M_{i}\) to \(M_{j}\) . Thus, all

\textit{[Image omitted: images/d924aac294cc9551361a30d6513d531c8f3272b538cfe80033e5e4c133bd5c92.jpg]}

flowchart

Policy management flowchart showing interactions between simulator,
rollout, policy, and update policy components

Figure 4: The task graph corresponding to an invocation of
train\_policy.remote() in Figure 3. Remote function calls and the actor
method calls correspond to tasks in the task graph. The figure shows two
actors. The method invocations for each actor (the tasks labeled
\(A_{1i}\) and \(A_{2i}\) ) have stateful edges between them indicating
that they share the mutable actor state. There are control edges from
train\_policy to the tasks that it invokes. To train multiple policies
in parallel, we could call train\_policy.remote() multiple times.

methods invoked on the same actor object form a chain that is connected
by stateful edges (Figure 4). This chain captures the order in which
these methods were invoked.

Stateful edges help us embed actors in an otherwise stateless task
graph, as they capture the implicit data dependency between successive
method invocations sharing the internal state of an actor. Stateful
edges also enable us to maintain lineage. As in other dataflow systems
{[}64{]}, we track data lineage to enable reconstruction. By explicitly
including stateful edges in the lineage graph, we can easily reconstruct
lost data, whether produced by remote functions or actor methods
(Section 4.2.3).

\section{4 Architecture}\label{architecture}

Ray's architecture comprises (1) an application layer implementing the
API, and (2) a system layer providing high scalability and fault
tolerance.

\section{4.1 Application Layer}\label{application-layer}

The application layer consists of three types of processes:

\begin{itemize}
\tightlist
\item
  Driver: A process executing the user program.\\
\item
  Worker: A stateless process that executes tasks (remote functions)
  invoked by a driver or another
\end{itemize}

\textit{[Image omitted: images/f3c36cc953eb1881f3b8237baaeae80e0af8816e2fc8cd96fedc66a519549b39.jpg]}

flowchart

\begin{Shaded}
\begin{Highlighting}[]
\NormalTok{graph TD}
\NormalTok{    subgraph App Layer}
\NormalTok{        A["Driver"] {-}{-}\textgreater{} B["Object Store"]}
\NormalTok{        C["Worker"] {-}{-}\textgreater{} B}
\NormalTok{        B {-}{-}\textgreater{} D["Local Scheduler"]}
\NormalTok{    end}
\NormalTok{    subgraph System Layer}
\NormalTok{        E["Global Scheduler"] {-}{-}\textgreater{} F["Global Control Store (GCS)"]}
\NormalTok{        F {-}{-}\textgreater{} G["Object Table"]}
\NormalTok{        F {-}{-}\textgreater{} H["Task Table"]}
\NormalTok{        F {-}{-}\textgreater{} I["Function Table"]}
\NormalTok{        F {-}{-}\textgreater{} J["Event Logs"]}
\NormalTok{    end}
\NormalTok{    subgraph Node Layer}
\NormalTok{        K["Actor"] {-}{-}\textgreater{} L["Object Store"]}
\NormalTok{        M["Driver"] {-}{-}\textgreater{} L}
\NormalTok{        L {-}{-}\textgreater{} N["Local Scheduler"]}
\NormalTok{    end}
\NormalTok{    subgraph Node Layer}
\NormalTok{        O["Worker"] {-}{-}\textgreater{} P["Object Store"]}
\NormalTok{        Q["Worker"] {-}{-}\textgreater{} P}
\NormalTok{        P {-}{-}\textgreater{} R["Local Scheduler"]}
\NormalTok{    end}
\NormalTok{    subgraph System Layer}
\NormalTok{        S["Web UI"] {-}{-}\textgreater{} T["Debugging Tools"]}
\NormalTok{        U["Profiling Tools"] {-}{-}\textgreater{} V["Error Diagnosis"]}
\NormalTok{    end}
\NormalTok{    A {-}{-}\textgreater{} K}
\NormalTok{    C {-}{-}\textgreater{} K}
\NormalTok{    B {-}{-}\textgreater{} K}
\NormalTok{    L {-}{-}\textgreater{} K}
\NormalTok{    N {-}{-}\textgreater{} K}
\NormalTok{    P {-}{-}\textgreater{} K}
\NormalTok{    R {-}{-}\textgreater{} K}
\NormalTok{    S {-}{-}\textgreater{} K}
\NormalTok{    T {-}{-}\textgreater{} K}
\NormalTok{    U {-}{-}\textgreater{} K}
\NormalTok{    V {-}{-}\textgreater{} K}
\NormalTok{    W["Error Diagnosis"] {-}{-}\textgreater{} T}
\NormalTok{    X["Error Diagnosis"] {-}{-}\textgreater{} U}
\NormalTok{    Y["Error Diagnosis"] {-}{-}\textgreater{} V}
\NormalTok{    Z["Error Diagnosis"] {-}{-}\textgreater{} W}
\end{Highlighting}
\end{Shaded}

Figure 5: Ray's architecture consists of two parts: an application layer
and a system layer. The application layer implements the API and the
computation model described in Section 3, the system layer implements
task scheduling and data management to satisfy the performance and
fault-tolerance requirements.

worker. Workers are started automatically and assigned tasks by the
system layer. When a remote function is declared, the function is
automatically published to all workers. A worker executes tasks
serially, with no local state maintained across tasks.

- Actor: A stateful process that executes, when invoked, only the
methods it exposes. Unlike a worker, an actor is explicitly instantiated
by a worker or a driver. Like workers, actors execute methods serially,
except that each method depends on the state resulting from the previous
method execution.

\section{4.2 System Layer}\label{system-layer}

The system layer consists of three major components: a global control
store, a distributed scheduler, and a distributed object store. All
components are horizontally scalable and fault-tolerant.

\section{4.2.1 Global Control Store
(GCS)}\label{global-control-store-gcs}

The global control store (GCS) maintains the entire control state of the
system, and it is a unique feature of our design. At its core, GCS is a
key-value store with pub-sub functionality. We use sharding to achieve
scale, and per-shard chain replication {[}61{]} to provide fault
tolerance. The primary reason for the GCS and its design is to maintain
fault tolerance and low latency for a system that can dynamically spawn
millions of tasks per second.

Fault tolerance in case of node failure requires a solution to maintain
lineage information. Existing lineage-based solutions {[}64, 63, 40,
28{]} focus on coarse-grained parallelism and can therefore use a single
node (e.g., master, driver) to store the lineage without impacting
performance. However, this design is not scalable for a fine-grained and
dynamic workload like simulation. Therefore, we decouple the durable
lineage storage from the other system components, allowing each to scale
independently.

Maintaining low latency requires minimizing overheads in task
scheduling, which involves choosing where to execute, and subsequently
task dispatch, which involves retrieving remote inputs from other nodes.
Many existing dataflow systems {[}64, 40, 48{]} couple these by storing
object locations and sizes in a centralized scheduler, a natural design
when the scheduler is not a bottleneck. However, the scale and
granularity that Ray targets requires keeping the centralized scheduler
off the critical path. Involving the scheduler in each object transfer
is prohibitively expensive for primitives important to distributed
training like allreduce, which is both communication-intensive and
latency-sensitive. Therefore, we store the object metadata in the GCS
rather than in the scheduler, fully decoupling task dispatch from task
scheduling.

In summary, the GCS significantly simplifies Ray's overall design, as it
enables every component in the system to be stateless. This not only
simplifies support for fault tolerance (i.e., on failure, components
simply restart and read the lineage from the GCS), but also makes it
easy to scale the distributed object store and scheduler independently,
as all components share the needed state via the GCS. An added benefit
is the easy development of debugging, profiling, and visualization
tools.

\section{4.2.2 Bottom-Up Distributed
Scheduler}\label{bottom-up-distributed-scheduler}

As discussed in Section 2, Ray needs to dynamically schedule millions of
tasks per second, tasks which may take as little as a few milliseconds.
None of the cluster schedulers we are aware of meet these requirements.
Most cluster computing frameworks, such as Spark {[}64{]}, CIEL
{[}40{]}, and Dryad {[}28{]} implement a centralized scheduler, which
can provide locality but at latencies in the tens of ms. Distributed
schedulers such as work stealing {[}12{]}, Sparrow {[}45{]} and Canary
{[}47{]} can achieve high scale, but they either don't consider data
locality {[}12{]}, or assume tasks belong to independent jobs {[}45{]},
or assume the computation graph is known {[}47{]}.

To satisfy the above requirements, we design a two-level hierarchical
scheduler consisting of a global scheduler and per-node local
schedulers. To avoid overloading the global scheduler, the tasks created
at a node are submitted first to the node's local scheduler. A local
scheduler schedules tasks locally unless the node is overloaded (i.e.,
its local task queue exceeds a predefined threshold), or it cannot
satisfy a task's requirements (e.g., lacks a GPU). If a local scheduler
decides not to schedule a task locally, it forwards it to the global
scheduler. Since this scheduler attempts to schedule tasks locally first
(i.e., at the leaves of the scheduling hierarchy), we call it a
bottom-up scheduler.

\textit{[Image omitted: images/084a3459ad3f2a84e1c531c8b7dbb02c39d674f8a6227bd335631725f2fb1d2b.jpg]}

flowchart

\begin{Shaded}
\begin{Highlighting}[]
\NormalTok{graph TD}
\NormalTok{    subgraph Node 1}
\NormalTok{        Driver {-}{-}\textgreater{} LocalScheduler}
\NormalTok{        Worker1 {-}{-}\textgreater{} LocalScheduler}
\NormalTok{        Worker2 {-}{-}\textgreater{} LocalScheduler}
\NormalTok{        WorkerN {-}{-}\textgreater{} LocalScheduler}
\NormalTok{    end}

\NormalTok{    subgraph Node N}
\NormalTok{        Worker1 {-}{-}\textgreater{} LocalScheduler}
\NormalTok{        Worker2 {-}{-}\textgreater{} LocalScheduler}
\NormalTok{        WorkerN {-}{-}\textgreater{} LocalScheduler}
\NormalTok{    end}

\NormalTok{    LocalScheduler {-}{-}\textgreater{} GlobalScheduler}
\NormalTok{    LocalScheduler {-}{-}\textgreater{} GlobalControlState}
\NormalTok{    GlobalControlState {-}{-}\textgreater{} GlobalScheduler}
\NormalTok{    GlobalScheduler {-}{-}\textgreater{} SubmitTasks}
\NormalTok{    GlobalControlState {-}{-}\textgreater{} ScheduleTasks}
\NormalTok{    GlobalControlState {-}{-}\textgreater{} LoadInfo}
\NormalTok{    GlobalControlState {-}{-}\textgreater{} Info}
\NormalTok{    style Node 1 fill:\#f9f,stroke:\#333}
\NormalTok{    style Node N fill:\#f9f,stroke:\#333}
\NormalTok{    style Node 1 fill:\#ccf,stroke:\#333}
\NormalTok{    style Node N fill:\#ccf,stroke:\#333}
\end{Highlighting}
\end{Shaded}

Figure 6: Bottom-up distributed scheduler. Tasks are submitted
bottom-up, from drivers and workers to a local scheduler and forwarded
to the global scheduler only if needed (Section 4.2.2). The thickness of
each arrow is proportional to its request rate.

The global scheduler considers each node's load and task's constraints
to make scheduling decisions. More precisely, the global scheduler
identifies the set of nodes that have enough resources of the type
requested by the task, and of these nodes selects the node which
provides the lowest estimated waiting time. At a given node, this time
is the sum of (i) the estimated time the task will be queued at that
node (i.e., task queue size times average task execution), and (ii) the
estimated transfer time of task's remote inputs (i.e., total size of
remote inputs divided by average bandwidth). The global scheduler gets
the queue size at each node and the node resource availability via
heartbeats, and the location of the task's inputs and their sizes from
GCS. Furthermore, the global scheduler computes the average task
execution and the average transfer bandwidth using simple exponential
averaging. If the global scheduler becomes a bottleneck, we can
instantiate more replicas all sharing the same information via GCS. This
makes our scheduler architecture highly scalable.

\section{4.2.3 In-Memory Distributed Object
Store}\label{in-memory-distributed-object-store}

To minimize task latency, we implement an in-memory distributed storage
system to store the inputs and outputs of every task, or stateless
computation. On each node, we implement the object store via shared
memory. This allows zero-copy data sharing between tasks running on the
same node. As a data format, we use Apache Arrow {[}2{]}.

If a task's inputs are not local, the inputs are replicated to the local
object store before execution. Also, a task writes its outputs to the
local object store. Replication eliminates the potential bottleneck due
to hot data objects and minimizes task execution time as a task only
reads/writes data from/to the local memory. This increases throughput
for computation-bound workloads, a profile shared by many AI
applications. For low latency, we keep objects entirely in memory and
evict them as needed to disk using an LRU policy.

As with existing cluster computing frameworks, such as Spark {[}64{]},
and Dryad {[}28{]}, the object store is limited to immutable data. This
obviates the need for complex consistency protocols (as objects are not
updated), and simplifies support for fault tolerance. In the case of
node failure, Ray recovers any needed objects through lineage
re-execution. The lineage stored in the GCS tracks both stateless tasks
and stateful actors during initial execution; we use the former to
reconstruct objects in the store.

For simplicity, our object store does not support distributed objects,
i.e., each object fits on a single node. Distributed objects like large
matrices or trees can be implemented at the application level as
collections of futures.

\section{4.2.4 Implementation}\label{implementation}

Ray is an active open source project \(^{†}\) developed at the
University of California, Berkeley. Ray fully integrates with the Python
environment and is easy to install by simply running pip install ray.
The implementation comprises \(\approx\) 40K lines of code (LoC), 72\%
in C++ for the system layer, 28\% in Python for the application layer.
The GCS uses one Redis {[}50{]} key-value store per shard, with entirely
single-key operations. GCS tables are sharded by object and task IDs to
scale, and every shard is chain-replicated {[}61{]} for fault tolerance.
We implement both the local and global schedulers as event-driven,
single-threaded processes. Internally, local schedulers maintain cached
state for local object metadata, tasks waiting for inputs, and tasks
ready for dispatch to a worker. To transfer large objects between
different object stores, we stripe the object across multiple TCP
connections.

\section{4.3 Putting Everything
Together}\label{putting-everything-together}

Figure 7 illustrates how Ray works end-to-end with a simple example that
adds two objects a and b, which could be scalars or matrices, and
returns result c.~The remote function add() is automatically registered
with the GCS upon initialization and distributed to every worker in the
system (step 0 in Figure 7a).

Figure 7a shows the step-by-step operations triggered by a driver
invoking add.remote(a, b), where a and b are stored on nodes N1 and N2,
respectively. The driver submits add(a, b) to the local scheduler (step
1), which forwards it to a global scheduler (step 2). \(^{\ddagger}\)
Next, the global scheduler looks up the locations of add(a, b)'s
arguments in the GCS (step 3) and decides to schedule the task on node
N2, which stores argument b (step 4). The local scheduler at node N2
checks whether the local object store contains add(a, b)'s arguments
(step 5). Since the

\textit{[Image omitted: images/bd2cc3e349470b131dabad825f431befe3931ed0b576f712a356618591f96410.jpg]}

flowchart

System architecture flowchart showing data flow between Driver, Function
Table, Worker, and Global Control Store (GCS) components with address
and return flags.

\begin{enumerate}
\def\labelenumi{(\alph{enumi})}
\tightlist
\item
  Executing a task remotely\\
  \textit{[Image omitted: images/46c79bc5e2f35e0ca7bb712bf579f9d25e605d32fe16aa55d1034acfab17a0c2.jpg]}
\end{enumerate}

flowchart

System architecture flowchart showing data flow between Driver, Global
Control Store (GCS), Function Table, Object Table, and Worker components
with labeled nodes and edges.

\begin{enumerate}
\def\labelenumi{(\alph{enumi})}
\setcounter{enumi}{1}
\tightlist
\item
  Returning the result of a remote task\\
  Figure 7: An end-to-end example that adds a and b and returns c.~Solid
  lines are data plane operations and dotted lines are control plane
  operations. (a) The function add() is registered with the GCS by node
  1 (N1), invoked on N1, and executed on N2. (b) N1 gets add()'s result
  using ray.get(). The Object Table entry for c is created in step 4 and
  updated in step 6 after c is copied to N1.
\end{enumerate}

local store doesn't have object \(a\) , it looks up \(a\) 's location in
the GCS (step 6). Learning that \(a\) is stored at \(N1\) , \(N2\) 's
object store replicates it locally (step 7). As all arguments of
\(\mathbf{add}()\) are now stored locally, the local scheduler invokes
\(\mathbf{add}()\) at a local worker (step 8), which accesses the
arguments via shared memory (step 9).

Figure 7b shows the step-by-step operations triggered by the execution
of ray.get() at N1, and of add() at N2, respectively. Upon
ray.get(idc)'s invocation, the driver checks the local object store for
the value c, using the future idc returned by add() (step 1). Since the
local object store doesn't store c, it looks up its location in the GCS.
At this time, there is no entry for c, as c has not been created yet. As
a result, N1's object store registers a callback with the Object Table
to be triggered when c's entry has been created (step 2). Meanwhile, at
N2, add() completes its execution, stores the result c in the local
object store (step 3), which in turn adds c's entry to the GCS (step 4).
As a result, the GCS triggers a callback to N1's object store with c's
entry (step 5). Next, N1 replicates c from N2 (step 6), and returns c to
ray.get() (step 7), which finally completes the task.

While this example involves a large number of RPCs, in many cases this
number is much smaller, as most tasks are scheduled locally, and the GCS
replies are cached by the global and local schedulers.

\textit{[Image omitted: images/8d715de7a0cdee4d60fee60c3aac3f87107b90df6dc63a5338bf1513784c492c.jpg]}

bar

{\def\LTcaptype{none} % do not increment counter
\begin{longtable}[]{@{}lll@{}}
\toprule\noalign{}
Object size & Locality Aware & Unaware \\
\midrule\noalign{}
\endhead
\bottomrule\noalign{}
\endlastfoot
100KB & 0.0001 & 0.0001 \\
1MB & 0.0001 & 0.0003 \\
10MB & 0.0001 & 0.003 \\
100MB & 0.0001 & 0.03 \\
\end{longtable}
}

\begin{enumerate}
\def\labelenumi{(\alph{enumi})}
\tightlist
\item
  Ray locality scheduling
\end{enumerate}

\textit{[Image omitted: images/9ec0de3a56c340def8de507a542f6e062afa6dddf396e14d94eb50a3832ac494.jpg]}

bar

{\def\LTcaptype{none} % do not increment counter
\begin{longtable}[]{@{}ll@{}}
\toprule\noalign{}
number of nodes & Millions of tasks/s \\
\midrule\noalign{}
\endhead
\bottomrule\noalign{}
\endlastfoot
10 & 0.1 \\
20 & 0.3 \\
30 & 0.5 \\
40 & 0.7 \\
50 & 0.9 \\
60 & 1.1 \\
100 & 1.6 \\
\end{longtable}
}

\begin{enumerate}
\def\labelenumi{(\alph{enumi})}
\setcounter{enumi}{1}
\tightlist
\item
  Ray scalability\\
  Figure 8: (a) Tasks leverage locality-aware placement. 1000 tasks with
  a random object dependency are scheduled onto one of two nodes. With
  locality-aware policy, task latency remains independent of the size of
  task inputs instead of growing by 1-2 orders of magnitude. (b)
  Near-linear scalability leveraging the GCS and bottom-up distributed
  scheduler. Ray reaches 1 million tasks per second throughput with 60
  nodes. \(x \in \{70, 80, 90\}\) omitted due to cost.
\end{enumerate}

\section{5 Evaluation}\label{evaluation}

In our evaluation, we study the following questions:

\begin{enumerate}
\def\labelenumi{\arabic{enumi}.}
\tightlist
\item
  How well does Ray meet the latency, scalability, and fault tolerance
  requirements listed in Section 2? (Section 5.1)\\
\item
  What overheads are imposed on distributed primitives (e.g., allreduce)
  written using Ray's API? (Section 5.1)\\
\item
  In the context of RL workloads, how does Ray compare against
  specialized systems for training, serving, and simulation? (Section
  5.2)\\
\item
  What advantages does Ray provide for RL applications, compared to
  custom systems? (Section 5.3)
\end{enumerate}

All experiments were run on Amazon Web Services. Unless otherwise
stated, we use m4.16xlarge CPU instances and p3.16xlarge GPU instances.

\section{5.1 Microbenchmarks}\label{microbenchmarks}

Locality-aware task placement. Fine-grain load balancing and
locality-aware placement are primary benefits of tasks in Ray. Actors,
once placed, are unable to move their computation to large remote
objects, while tasks can. In Figure 8a, tasks placed without data
locality awareness (as is the case for actor methods), suffer 1-2 orders
of magnitude latency increase at 10-100MB input data sizes. Ray unifies
tasks and actors through the shared object store, allowing developers to
use tasks for e.g., expensive postprocessing on output produced by
simulation actors.

End-to-end scalability. One of the key benefits of

\textit{[Image omitted: images/b4c1baaed8e2575824aa8678457645a5ed71af5d77ffba2fbe9b19f5782f5d9a.jpg]}

bar\_line

{\def\LTcaptype{none} % do not increment counter
\begin{longtable}[]{@{}lll@{}}
\toprule\noalign{}
object size & IOPS & throughput (GB/s) \\
\midrule\noalign{}
\endhead
\bottomrule\noalign{}
\endlastfoot
1KB & 18000 & 0 \\
10KB & 17000 & 0 \\
100KB & 11000 & 1 \\
1MB & 2000 & 2 \\
10MB & 11000 & 9 \\
100MB & 17000 & 14 \\
1GB & 18000 & 15 \\
\end{longtable}
}

Figure 9: Object store write throughput and IOPS. From a single client,
throughput exceeds 15GB/s (red) for large objects and 18K IOPS (cyan)
for small objects on a 16 core instance (m4.4xlarge). It uses 8 threads
to copy objects larger than 0.5MB and 1 thread for small objects. Bar
plots report throughput with 1, 2, 4, 8, 16 threads. Results are
averaged over 5 runs.

the Global Control Store (GCS) and the bottom-up distributed scheduler
is the ability to horizontally scale the system to support a high
throughput of fine-grained tasks, while maintaining fault tolerance and
low-latency task scheduling. In Figure 8b, we evaluate this ability on
an embarrassingly parallel workload of empty tasks, increasing the
cluster size on the x-axis. We observe near-perfect linearity in
progressively increasing task throughput. Ray exceeds 1 million tasks
per second throughput at 60 nodes and continues to scale linearly beyond
1.8 million tasks per second at 100 nodes. The rightmost datapoint shows
that Ray can process 100 million tasks in less than a minute (54s), with
minimum variability. As expected, increasing task duration reduces
throughput proportionally to mean task duration, but the overall
scalability remains linear. While many realistic workloads may exhibit
more limited scalability due to object dependencies and inherent limits
to application parallelism, this demonstrates the scalability of our
overall architecture under high load.

Object store performance. To evaluate the performance of the object
store (Section 4.2.3), we track two metrics: IOPS (for small objects)
and write throughput (for large objects). In Figure 9, the write
throughput from a single client exceeds 15GB/s as object size increases.
For larger objects, memcpy dominates object creation time. For smaller
objects, the main overheads are in serialization and IPC between the
client and object store.

GCS fault tolerance. To maintain low latency while providing strong
consistency and fault tolerance, we build a lightweight chain
replication {[}61{]} layer on top of Redis. Figure 10a simulates
recording Ray tasks to and reading tasks from the GCS, where keys are 25
bytes and values are 512 bytes. The client sends requests as fast as it
can, having at most one in-flight request at a time. Failures are
reported to the chain master either from the client (having received
explicit errors, or timeouts despite retries) or from any server in the
chain (having received explicit errors). Overall, reconfigurations
caused a maximum client-observed delay of under 30ms (this includes both
failure detection and recovery delays).

\textit{[Image omitted: images/ec3f6dcdbeb402d883a8933af3ba5d502bd4e89639114da28288d7aa88e6cdf9.jpg]}

line

{\def\LTcaptype{none} % do not increment counter
\begin{longtable}[]{@{}llll@{}}
\toprule\noalign{}
Time since start (s) & write (μs) & read (μs) & node dead (μs) \\
\midrule\noalign{}
\endhead
\bottomrule\noalign{}
\endlastfoot
0 & \textasciitilde500 & \textasciitilde200 & - \\
1 & \textasciitilde500 & \textasciitilde200 & - \\
2 & \textasciitilde500 & \textasciitilde200 & - \\
3 & \textasciitilde500 & \textasciitilde200 & - \\
4 & \textasciitilde500 & \textasciitilde200 & - \\
5 & \textasciitilde10000 & \textasciitilde10000 &
\textasciitilde10000 \\
6 & \textasciitilde500 & \textasciitilde200 & - \\
7 & \textasciitilde500 & \textasciitilde200 & - \\
8 & \textasciitilde500 & \textasciitilde200 & - \\
9 & \textasciitilde500 & \textasciitilde200 & - \\
10 & \textasciitilde500 & \textasciitilde200 & - \\
\end{longtable}
}

\begin{enumerate}
\def\labelenumi{(\alph{enumi})}
\tightlist
\item
  A timeline for GCS read and write latencies as viewed from a client
  submitting tasks. The chain starts with 2 replicas. We manually
  trigger reconfiguration as follows. At \(t \approx 4.2s\) , a chain
  member is killed; immediately after, a new chain member joins,
  initiates state transfer, and restores the chain to 2-way replication.
  The maximum client-observed latency is under 30ms despite
  reconfigurations.\\
  \textit{[Image omitted: images/c64994a936c32a39d8b691257ba5bbcd4ad62ef55f206484e59de2a132cbdb05.jpg]}
\end{enumerate}

line

{\def\LTcaptype{none} % do not increment counter
\begin{longtable}[]{@{}
  >{\raggedright\arraybackslash}p{(\linewidth - 4\tabcolsep) * \real{0.3492}}
  >{\raggedright\arraybackslash}p{(\linewidth - 4\tabcolsep) * \real{0.3492}}
  >{\raggedright\arraybackslash}p{(\linewidth - 4\tabcolsep) * \real{0.3016}}@{}}
\toprule\noalign{}
\begin{minipage}[b]{\linewidth}\raggedright
Elapsed Time (seconds)
\end{minipage} & \begin{minipage}[b]{\linewidth}\raggedright
Ray, no GCS flush (MB)
\end{minipage} & \begin{minipage}[b]{\linewidth}\raggedright
Ray, GCS flush (MB)
\end{minipage} \\
\midrule\noalign{}
\endhead
\bottomrule\noalign{}
\endlastfoot
0 & 0 & 0 \\
5000 & 8000 & 0 \\
60000 & 8000 & 0 \\
\end{longtable}
}

\begin{enumerate}
\def\labelenumi{(\alph{enumi})}
\setcounter{enumi}{1}
\tightlist
\item
  The Ray GCS maintains a constant memory footprint with GCS flushing.
  Without GCS flushing, the memory footprint reaches a maximum capacity
  and the workload fails to complete within a predetermined duration
  (indicated by the red cross).\\
  Figure 10: Ray GCS fault tolerance and flushing.
\end{enumerate}

GCS flushing. Ray is equipped to periodically flush the contents of GCS
to disk. In Figure 10b we submit 50 million empty tasks sequentially and
monitor GCS memory consumption. As expected, it grows linearly with the
number of tasks tracked and eventually reaches the memory capacity of
the system. At that point, the system becomes stalled and the workload
fails to finish within a reasonable amount of time. With periodic GCS
flushing, we achieve two goals. First, the memory footprint is capped at
a user-configurable level (in the microbenchmark we employ an aggressive
strategy where consumed memory is kept as low as possible). Second, the
flushing mechanism provides a natural way to snapshot lineage to disk
for long-running Ray applications.

Recovering from task failures. In Figure 11a, we

\textit{[Image omitted: images/bf23b95e69314f66f2108f7707a9f6045a4bf38b6962cc62720d82f25ce2a081.jpg]}

bar\_line

{\def\LTcaptype{none} % do not increment counter
\begin{longtable}[]{@{}lll@{}}
\toprule\noalign{}
Time since start (s) & Throughput (tasks/s) & Number of nodes \\
\midrule\noalign{}
\endhead
\bottomrule\noalign{}
\endlastfoot
0 & 2000 & 60 \\
50 & 1000 & 20 \\
100 & 500 & 10 \\
150 & 500 & 10 \\
200 & 500 & 10 \\
250 & 2000 & 60 \\
\end{longtable}
}

\begin{enumerate}
\def\labelenumi{(\alph{enumi})}
\tightlist
\item
  Task reconstruction\\
  \textit{[Image omitted: images/aaeb2a17cae1576dac2a05cb169142551df6c9fe6df37a6ea0995e1654cfe4e2.jpg]}
\end{enumerate}

line

{\def\LTcaptype{none} % do not increment counter
\begin{longtable}[]{@{}
  >{\raggedright\arraybackslash}p{(\linewidth - 6\tabcolsep) * \real{0.2985}}
  >{\raggedright\arraybackslash}p{(\linewidth - 6\tabcolsep) * \real{0.2090}}
  >{\raggedright\arraybackslash}p{(\linewidth - 6\tabcolsep) * \real{0.2537}}
  >{\raggedright\arraybackslash}p{(\linewidth - 6\tabcolsep) * \real{0.2388}}@{}}
\toprule\noalign{}
\begin{minipage}[b]{\linewidth}\raggedright
Time since start (s)
\end{minipage} & \begin{minipage}[b]{\linewidth}\raggedright
Original tasks
\end{minipage} & \begin{minipage}[b]{\linewidth}\raggedright
Re-executed tasks
\end{minipage} & \begin{minipage}[b]{\linewidth}\raggedright
Checkpoint tasks
\end{minipage} \\
\midrule\noalign{}
\endhead
\bottomrule\noalign{}
\endlastfoot
0 & 400 & 0 & 400 \\
100 & 400 & 0 & 400 \\
200 & 400 & 0 & 400 \\
300 & 300 & 0 & 300 \\
400 & 300 & 0 & 300 \\
500 & 300 & 0 & 300 \\
600 & 300 & 0 & 300 \\
\end{longtable}
}

\begin{enumerate}
\def\labelenumi{(\alph{enumi})}
\setcounter{enumi}{1}
\tightlist
\item
  Actor reconstruction\\
  Figure 11: Ray fault-tolerance. (a) Ray reconstructs lost task
  dependencies as nodes are removed (dotted line), and recovers to
  original throughput when nodes are added back. Each task is 100ms and
  depends on an object generated by a previously submitted task. (b)
  Actors are reconstructed from their last checkpoint. At t = 200s, we
  kill 2 of the 10 nodes, causing 400 of the 2000 actors in the cluster
  to be recovered on the remaining nodes (t = 200--270s).
\end{enumerate}

demonstrate Ray's ability to transparently recover from worker node
failures and elastically scale, using the durable GCS lineage storage.
The workload, run on m4.xlarge instances, consists of linear chains of
100ms tasks submitted by the driver. As nodes are removed (at 25s, 50s,
100s), the local schedulers reconstruct previous results in the chain in
order to continue execution. Overall per-node throughput remains stable
throughout.

Recovering from actor failures. By encoding actor method calls as
stateful edges directly in the dependency graph, we can reuse the same
object reconstruction mechanism as in Figure 11a to provide transparent
fault tolerance for stateful computation. Ray additionally leverages
user-defined checkpoint functions to bound the reconstruction time for
actors (Figure 11b). With minimal overhead, checkpointing enables only
500 methods to be re-executed, versus 10k re-executions without
checkpointing. In the future, we hope to further reduce actor
reconstruction time, e.g., by allowing users to annotate methods that do
not mutate state.

Allreduce. Allreduce is a distributed communication primitive important
to many machine learning workloads. Here, we evaluate whether Ray can
natively support a ring allreduce {[}57{]} implementation with low
enough overhead to match existing implementations {[}53{]}. We find that
Ray completes allreduce across 16 nodes on 100MB in \textasciitilde200ms
and 1GB in \textasciitilde1200ms, surprisingly outperforming OpenMPI
(v1.10), a popular MPI implementation, by 1.5× and 2× respectively
(Figure 12a). We attribute Ray's performance to its use of multiple
threads for network transfers, taking full advantage of the 25Gbps
connection between nodes on AWS, whereas OpenMPI sequentially sends and
receives data on a single thread {[}22{]}. For smaller objects, OpenMPI
outperforms Ray by switching to a lower overhead algorithm, an
optimization we plan to implement in the future.

\textit{[Image omitted: images/bb5eeef50f37ef0209e03dee567d6fcd55ea9584718b18895fa122db0ae69a56.jpg]}

bar

{\def\LTcaptype{none} % do not increment counter
\begin{longtable}[]{@{}llll@{}}
\toprule\noalign{}
Object size & OpenMPI & Ray* & Ray \\
\midrule\noalign{}
\endhead
\bottomrule\noalign{}
\endlastfoot
10MB & 30 & 100 & 100 \\
100MB & 300 & 200 & 150 \\
1GB & 3000 & 2000 & 1000 \\
\end{longtable}
}

\begin{enumerate}
\def\labelenumi{(\alph{enumi})}
\tightlist
\item
  Ray vs OpenMPI
\end{enumerate}

\textit{[Image omitted: images/a29f97552f5e0a1c0ee35c01bf80d43dbc26fa73f09811c0f889d0e958f79ae0.jpg]}

bar

{\def\LTcaptype{none} % do not increment counter
\begin{longtable}[]{@{}ll@{}}
\toprule\noalign{}
Added scheduler latency (ms) & Iteration time (milliseconds) \\
\midrule\noalign{}
\endhead
\bottomrule\noalign{}
\endlastfoot
+0 & 200 \\
+1 & 250 \\
+5 & 350 \\
+10 & 550 \\
\end{longtable}
}

\begin{enumerate}
\def\labelenumi{(\alph{enumi})}
\setcounter{enumi}{1}
\tightlist
\item
  Ray scheduler ablation\\
  Figure 12: (a) Mean execution time of allreduce on 16 m4.16xl nodes.
  Each worker runs on a distinct node. Ray* restricts Ray to 1 thread
  for sending and 1 thread for receiving. (b) Ray's low-latency
  scheduling is critical for allreduce.
\end{enumerate}

Ray's scheduler performance is critical to implementing primitives such
as allreduce. In Figure 12b, we inject artificial task execution delays
and show that performance drops nearly \(2 \times\) with just a few ms
of extra latency. Systems with centralized schedulers like Spark and
CIEL typically have scheduler overheads in the tens of milliseconds
{[}62, 38{]}, making such workloads impractical. Scheduler throughput
also becomes a bottleneck since the number of tasks required by ring
reduce scales quadratically with the number of participants.

\section{5.2 Building blocks}\label{building-blocks}

End-to-end applications (e.g., AlphaGo {[}54{]}) require a tight
coupling of training, serving, and simulation. In this section, we
isolate each of these workloads to a setting that illustrates a typical
RL application's requirements. Due to a flexible programming model
targeted to RL, and a system designed to support this programming model,
Ray matches and sometimes exceeds the performance of dedicated systems
for these individual workloads.

\textit{[Image omitted: images/42fe127874ee7f037a85071da592d2711d035203f384803c49d5c3a80acd8a2f.jpg]}

bar

{\def\LTcaptype{none} % do not increment counter
\begin{longtable}[]{@{}llll@{}}
\toprule\noalign{}
Num GPUs (V100) & Horovod + TF & Distributed TF & Ray + TF \\
\midrule\noalign{}
\endhead
\bottomrule\noalign{}
\endlastfoot
4 & 800 & 800 & 800 \\
8 & 1600 & 1600 & 1600 \\
16 & 1600 & 2000 & 1800 \\
32 & 3000 & 3800 & 3200 \\
64 & 6000 & 6800 & 6000 \\
\end{longtable}
}

Figure 13: Images per second reached when distributing the training of a
ResNet-101 TensorFlow model (from the official TF benchmark). All
experiments were run on p3.16xl instances connected by 25Gbps Ethernet,
and workers allocated 4 GPUs per node as done in Horovod {[}53{]}. We
note some measurement deviations from previously reported, likely due to
hardware differences and recent TensorFlow performance improvements. We
used OpenMPI 3.0, TF 1.8, and NCCL2 for all runs.

\section{5.2.1 Distributed Training}\label{distributed-training}

We implement data-parallel synchronous SGD leveraging the Ray actor
abstraction to represent model replicas. Model weights are synchronized
via allreduce (5.1) or parameter server, both implemented on top of the
Ray API.

In Figure 13, we evaluate the performance of the Ray (synchronous)
parameter-server SGD implementation against state-of-the-art
implementations {[}53{]}, using the same TensorFlow model and synthetic
data generator for each experiment. We compare only against
TensorFlow-based systems to accurately measure the overhead imposed by
Ray, rather than differences between the deep learning frameworks
themselves. In each iteration, model replica actors compute gradients in
parallel, send the gradients to a sharded parameter server, then read
the summed gradients from the parameter server for the next iteration.

Figure 13 shows that Ray matches the performance of Horovod and is
within 10\% of distributed TensorFlow (in distributed\_replicated mode).
This is due to the ability to express the same application-level
optimizations found in these specialized systems in Ray's
general-purpose API. A key optimization is the pipelining of gradient
computation, transfer, and summation within a single iteration. To
overlap GPU computation with network transfer, we use a custom
TensorFlow operator to write tensors directly to Ray's object store.

\section{5.2.2 Serving}\label{serving}

Model serving is an important component of end-to-end applications. Ray
focuses primarily on the embedded serving of models to simulators
running within the same dynamic task graph (e.g., within an RL
application on Ray). In contrast, systems like Clipper {[}19{]} focus on
serving predictions to external clients.

In this setting, low latency is critical for achieving high utilization.
To show this, in Table 3 we compare the server throughput achieved using
a Ray actor to serve a policy versus using the open source Clipper
system over REST. Here, both client and server processes are colocated
on the same machine (a p3.8xlarge instance). This is often the case for
RL applications but not for the general web serving workloads addressed
by systems like Clipper. Due to its low-overhead serialization and
shared memory abstractions, Ray achieves an order of magnitude higher
throughput for a small fully connected policy model that takes in a
large input and is also faster on a more expensive residual network
policy model, similar to one used in AlphaGo Zero, that takes smaller
input.

System

Small Input

Larger Input

Clipper

4400 ± 15 states/sec

290 ± 1.3 states/sec

Ray

6200 ± 21 states/sec

6900 ± 150 states/sec

Table 3: Throughput comparisons for Clipper {[}19{]}, a dedicated
serving system, and Ray for two embedded serving workloads. We use a
residual network and a small fully connected network, taking 10ms and
5ms to evaluate, respectively. The server is queried by clients that
each send states of size 4KB and 100KB respectively in batches of 64.

\section{5.2.3 Simulation}\label{simulation}

Simulators used in RL produce results with variable lengths
(``timesteps'') that, due to the tight loop with training, must be used
as soon as they are available. The task heterogeneity and timeliness
requirements make simulations hard to support efficiently in BSP-style
systems. To demonstrate, we compare (1) an MPI implementation that
submits 3n parallel simulation runs on n cores in 3 rounds, with a
global barrier between rounds \(^{§}\) , to (2) a Ray program that
issues the same 3n tasks while concurrently gathering simulation results
back to the driver. Table 4 shows that both systems scale well, yet Ray
achieves up to \(1.8\times\) throughput. This motivates a programming
model that can dynamically spawn and collect the results of fine-grained
simulation tasks.

System, programming model

1 CPU

16 CPUs

256 CPUs

MPI, bulk synchronous

22.6K

208K

2.16M

Ray, asynchronous tasks

22.3K

290K

4.03M

Table 4: Timesteps per second for the Pendulum-v0 simulator in OpenAI
Gym {[}13{]}. Ray allows for better utilization when running
heterogeneous simulations at scale.

\section{5.3 RL Applications}\label{rl-applications}

Without a system that can tightly couple the training, simulation, and
serving steps, reinforcement learning algorithms today are implemented
as one-off solutions that make it difficult to incorporate optimizations
that, for example, require a different computation structure or that
utilize different architectures. Consequently, with implementations of
two representative reinforcement learning applications in Ray, we are
able to match and even outperform custom systems built specifically for
these algorithms. The primary reason is the flexibility of Ray's
programming model, which can express application-level optimizations
that would require substantial engineering effort to port to
custom-built systems, but are transparently supported by Ray's dynamic
task graph execution engine.

\section{5.3.1 Evolution Strategies}\label{evolution-strategies}

To evaluate Ray on large-scale RL workloads, we implement the evolution
strategies (ES) algorithm and compare to the reference implementation
{[}49{]}---a system specially built for this algorithm that relies on
Redis for messaging and low-level multiprocessing libraries for
data-sharing. The algorithm periodically broadcasts a new policy to a
pool of workers and aggregates the results of roughly 10000 tasks (each
performing 10 to 1000 simulation steps).

As shown in Figure 14a, an implementation on Ray scales to 8192 cores.
Doubling the cores available yields an average completion time speedup
of 1.6×. Conversely, the special-purpose system fails to complete at
2048 cores, where the work in the system exceeds the processing capacity
of the application driver. To avoid this issue, the Ray implementation
uses an aggregation tree of actors, reaching a median time of 3.7
minutes, more than twice as fast as the best published result (10
minutes).

Initial parallelization of a serial implementation using Ray required
modifying only 7 lines of code. Performance improvement through
hierarchical aggregation was easy to realize with Ray's support for
nested tasks and actors. In contrast, the reference implementation had
several hundred lines of code dedicated to a protocol for communicating
tasks and data between workers, and would require further engineering to
support optimizations like hierarchical aggregation.

\section{5.3.2 Proximal Policy
Optimization}\label{proximal-policy-optimization}

We implement Proximal Policy Optimization (PPO) {[}51{]} in Ray and
compare to a highly-optimized reference implementation {[}5{]} that uses
OpenMPI communication primitives. The algorithm is an asynchronous
scatter-gather, where new tasks are assigned to simulation actors as
they return rollouts to the driver. Tasks are submitted until 320000
simulation steps are collected (each task produces between 10 and 1000
steps). The policy update performs 20 steps of SGD with a batch size of
32768. The model parameters in this example are roughly 350KB. These
experiments were run using p2.16xlarge (GPU) and m4.16xlarge (high CPU)
instances.

\textit{[Image omitted: images/d9869482612dddd7bda9dd5f94b813172369ce21fc7eeb4a22125627e3c2d277.jpg]}

bar

{\def\LTcaptype{none} % do not increment counter
\begin{longtable}[]{@{}lll@{}}
\toprule\noalign{}
Number of CPUs & Reference ES & Ray ES \\
\midrule\noalign{}
\endhead
\bottomrule\noalign{}
\endlastfoot
256 & 65 & 40 \\
1024 & 35 & 25 \\
8192 & 15 & 10 \\
\end{longtable}
}

\begin{enumerate}
\def\labelenumi{(\alph{enumi})}
\tightlist
\item
  Evolution Strategies
\end{enumerate}

\textit{[Image omitted: images/e2244a7933723590dd66fd899ec1e754cab007c09ac180fae93c7112f03802bb.jpg]}

bar

{\def\LTcaptype{none} % do not increment counter
\begin{longtable}[]{@{}lll@{}}
\toprule\noalign{}
CPUs x GPUs & MPI PPO & Ray PPO \\
\midrule\noalign{}
\endhead
\bottomrule\noalign{}
\endlastfoot
8x1 & 500 & 380 \\
64x8 & 150 & 100 \\
512x64 & 50 & 30 \\
\end{longtable}
}

\begin{enumerate}
\def\labelenumi{(\alph{enumi})}
\setcounter{enumi}{1}
\tightlist
\item
  PPO\\
  Figure 14: Time to reach a score of 6000 in the Humanoid-v1 task
  {[}13{]}. (a) The Ray ES implementation scales well to 8192 cores and
  achieves a median time of 3.7 minutes, over twice as fast as the best
  published result. The special-purpose system failed to run beyond 1024
  cores. ES is faster than PPO on this benchmark, but shows greater
  runtime variance. (b) The Ray PPO implementation outperforms a
  specialized MPI implementation {[}5{]} with fewer GPUs, at a fraction
  of the cost. The MPI implementation required 1 GPU for every 8 CPUs,
  whereas the Ray version required at most 8 GPUs (and never more than 1
  GPU per 8 CPUs).
\end{enumerate}

As shown in Figure 14b, the Ray implementation outperforms the optimized
MPI implementation in all experiments, while using a fraction of the
GPUs. The reason is that Ray is heterogeneity-aware and allows the user
to utilize asymmetric architectures by expressing resource requirements
at the granularity of a task or actor. The Ray implementation can then
leverage TensorFlow's single-process multi-GPU support and can pin
objects in GPU memory when possible. This optimization cannot be easily
ported to MPI due to the need to asynchronously gather rollouts to a
single GPU process. Indeed, {[}5{]} includes two custom implementations
of PPO, one using MPI for large clusters and one that is optimized for
GPUs but that is restricted to a single node. Ray allows for an
implementation suitable for both scenarios.

Ray's ability to handle resource heterogeneity also decreased PPO's cost
by a factor of 4.5 {[}4{]}, since CPU-only tasks can be scheduled on
cheaper high-CPU instances. In contrast, MPI applications often exhibit
symmetric architectures, in which all processes run the same code and
require identical resources, in this case preventing the use of CPU-only
machines for scale-out. Furthermore, the MPI implementation requires
on-demand instances since it does not transparently handle failure.
Assuming \(4 \times\) cheaper spot instances, Ray's fault tolerance and
resource-aware scheduling together cut costs by \(18 \times\) .

\section{6 Related Work}\label{related-work}

Dynamic task graphs. Ray is closely related to CIEL {[}40{]} and Dask
{[}48{]}. All three support dynamic task graphs with nested tasks and
implement the futures abstraction. CIEL also provides lineage-based
fault tolerance, while Dask, like Ray, fully integrates with Python.
However, Ray differs in two aspects that have important performance
consequences. First, Ray extends the task model with an actor
abstraction. This is necessary for efficient stateful computation in
distributed training and serving, to keep the model data collocated with
the computation. Second, Ray employs a fully distributed and decoupled
control plane and scheduler, instead of relying on a single master
storing all metadata. This is critical for efficiently supporting
primitives like allreduce without system modification. At peak
performance for 100MB on 16 nodes, allreduce on Ray (Section 5.1)
submits 32 rounds of 16 tasks in 200ms. Meanwhile, Dask reports a
maximum scheduler throughput of 3k tasks/s on 512 cores {[}3{]}. With a
centralized scheduler, each round of allreduce would then incur a
minimum of \(\sim\) 5ms of scheduling delay, translating to up to 2×
worse completion time (Figure 12b). Even with a decentralized scheduler,
coupling the control plane information with the scheduler leaves the
latter on the critical path for data transfer, adding an extra roundtrip
to every round of allreduce.

Dataflow systems. Popular dataflow systems, such as MapReduce {[}20{]},
Spark {[}65{]}, and Dryad {[}28{]} have widespread adoption for
analytics and ML workloads, but their computation model is too
restrictive for a fine-grained and dynamic simulation workload. Spark
and MapReduce implement the BSP execution model, which assumes that
tasks within the same stage perform the same computation and take
roughly the same amount of time. Dryad relaxes this restriction but
lacks support for dynamic task graphs. Furthermore, none of these
systems provide an actor abstraction, nor implement a distributed
scalable control plane and scheduler. Finally, Naiad {[}39{]} is a
dataflow system that provides improved scalability for some workloads,
but only supports static task graphs.

Machine learning frameworks. TensorFlow {[}7{]} and MXNet {[}18{]}
target deep learning workloads and efficiently leverage both CPUs and
GPUs. While they achieve great performance for training workloads
consisting of static DAGs of linear algebra operations, they have
limited support for the more general computation required to tightly
couple training with simulation and embedded serving. TensorFlow Fold
{[}33{]} provides some support for dynamic task graphs, as well as MXNet
through its internal C++ APIs, but neither fully supports the ability to
modify the DAG during execution in response to task progress, task
completion times, or faults. TensorFlow and MXNet in principle achieve
generality by allowing the programmer to simulate low-level
message-passing and synchronization primitives, but the pitfalls and
user experience in this case are similar to those of MPI. OpenMPI
{[}22{]} can achieve high performance, but it is relatively hard to
program as it requires explicit coordination to handle heterogeneous and
dynamic task graphs. Furthermore, it forces the programmer to explicitly
handle fault tolerance.

Actor systems. Orleans {[}14{]} and Akka {[}1{]} are two actor
frameworks well suited to developing highly available and concurrent
distributed systems. However, compared to Ray, they provide less support
for recovery from data loss. To recover stateful actors, the Orleans
developer must explicitly checkpoint actor state and intermediate
responses. Stateless actors in Orleans can be replicated for scale-out,
and could therefore act as tasks, but unlike in Ray, they have no
lineage. Similarly, while Akka explicitly supports persisting actor
state across failures, it does not provide efficient fault tolerance for
stateless computation (i.e., tasks). For message delivery, Orleans
provides at-least-once and Akka provides at-most-once semantics. In
contrast, Ray provides transparent fault tolerance and exactly-once
semantics, as each method call is logged in the GCS and both arguments
and results are immutable. We find that in practice these limitations do
not affect the performance of our applications. Erlang {[}10{]} and C++
Actor Framework {[}17{]}, two other actor-based systems, have similarly
limited support for fault tolerance.

Global control store and scheduling. The concept of logically
centralizing the control plane has been previously proposed in software
defined networks (SDNs) {[}16{]}, distributed file systems (e.g., GFS
{[}23{]}), resource management (e.g., Omega {[}52{]}), and distributed
frameworks (e.g., MapReduce {[}20{]}, BOOM {[}9{]}), to name a few. Ray
draws inspiration from these pioneering efforts, but provides
significant improvements. In contrast with SDNs, BOOM, and GFS, Ray
decouples the storage of the control plane information (e.g., GCS) from
the logic implementation (e.g., schedulers). This allows both storage
and computation layers to scale independently, which is key to achieving
our scalability targets. Omega uses a distributed architecture in which
schedulers coordinate via globally shared state. To this architecture,
Ray adds global schedulers to balance load across local schedulers, and
targets ms-level, not second-level, task scheduling.

Ray implements a unique distributed bottom-up scheduler that is
horizontally scalable, and can handle dynamically constructed task
graphs. Unlike Ray, most existing cluster computing systems {[}20, 64,
40{]} use a centralized scheduler architecture. While Sparrow {[}45{]}
is decentralized, its schedulers make independent decisions, limiting
the possible scheduling policies, and all tasks of a job are handled by
the same global scheduler. Mesos {[}26{]} implements a two-level
hierarchical scheduler, but its top-level scheduler manages frameworks,
not individual tasks.

Canary {[}47{]} achieves impressive performance by having each scheduler
instance handle a portion of the task graph, but does not handle dynamic
computation graphs.

Cilk {[}12{]} is a parallel programming language whose work-stealing
scheduler achieves provably efficient load-balancing for dynamic task
graphs. However, with no central coordinator like Ray's global
scheduler, this fully parallel design is also difficult to extend to
support data locality and resource heterogeneity in a distributed
setting.

\section{7 Discussion and Experiences}\label{discussion-and-experiences}

Building Ray has been a long journey. It started two years ago with a
Spark library to perform distributed training and simulations. However,
the relative inflexibility of the BSP model, the high per-task overhead,
and the lack of an actor abstraction led us to develop a new system.
Since we released Ray roughly one year ago, several hundreds of people
have used it and several companies are running it in production. Here we
discuss our experience developing and using Ray, and some early user
feedback.

API. In designing the API, we have emphasized minimalism. Initially we
started with a basic task abstraction. Later, we added the wait()
primitive to accommodate rollouts with heterogeneous durations and the
actor abstraction to accommodate third-party simulators and amortize the
overhead of expensive initializations. While the resulting API is
relatively low-level, it has proven both powerful and simple to use. We
have already used this API to implement many state-of-the-art RL
algorithms on top of Ray, including A3C {[}36{]}, PPO {[}51{]}, DQN
{[}37{]}, ES {[}49{]}, DDPG {[}55{]}, and Ape-X {[}27{]}. In most cases
it took us just a few tens of lines of code to port these algorithms to
Ray. Based on early user feedback, we are considering enhancing the API
to include higher level primitives and libraries, which could also
inform scheduling decisions.

Limitations. Given the workload generality, specialized optimizations
are hard. For example, we must make scheduling decisions without full
knowledge of the computation graph. Scheduling optimizations in Ray
might require more complex runtime profiling. In addition, storing
lineage for each task requires the implementation of garbage collection
policies to bound storage costs in the GCS, a feature we are actively
developing.

Fault tolerance. We are often asked if fault tolerance is really needed
for AI applications. After all, due to the statistical nature of many AI
algorithms, one could simply ignore failed rollouts. Based on our
experience, our answer is ``yes''. First, the ability to ignore failures
makes applications much easier to write and reason about. Second, our
particular implementation of fault tolerance via deterministic replay
dramatically simplifies debugging as it allows us to easily reproduce
most errors. This is particularly important since, due to their
stochasticity, AI algorithms are notoriously hard to debug. Third, fault
tolerance helps save money since it allows us to run on cheap resources
like spot instances on AWS. Of course, this comes at the price of some
overhead. However, we found this overhead to be minimal for our target
workloads.

GCS and Horizontal Scalability. The GCS dramatically simplified Ray
development and debugging. It enabled us to query the entire system
state while debugging Ray itself, instead of having to manually expose
internal component state. In addition, the GCS is also the backend for
our timeline visualization tool, used for application-level debugging.

The GCS was also instrumental to Ray's horizontal scalability. In
Section 5, we were able to scale by adding more shards whenever the GCS
became a bottleneck. The GCS also enabled the global scheduler to scale
by simply adding more replicas. Due to these advantages, we believe that
centralizing control state will be a key design component of future
distributed systems.

\section{8 Conclusion}\label{conclusion}

No general-purpose system today can efficiently support the tight loop
of training, serving, and simulation. To express these core building
blocks and meet the demands of emerging AI applications, Ray unifies
task-parallel and actor programming models in a single dynamic task
graph and employs a scalable architecture enabled by the global control
store and a bottom-up distributed scheduler. The programming
flexibility, high throughput, and low latencies simultaneously achieved
by this architecture is particularly important for emerging artificial
intelligence workloads, which produce tasks diverse in their resource
requirements, duration, and functionality. Our evaluation demonstrates
linear scalability up to 1.8 million tasks per second, transparent fault
tolerance, and substantial performance improvements on several
contemporary RL workloads. Thus, Ray provides a powerful combination of
flexibility, performance, and ease of use for the development of future
AI applications.

\section{9 Acknowledgments}\label{acknowledgments}

This research is supported in part by NSF CISE Expeditions Award
CCF-1730628 and gifts from Alibaba, Amazon Web Services, Ant Financial,
Arm, CapitalOne, Ericsson, Facebook, Google, Huawei, Intel, Microsoft,
Scotiabank, Splunk and VMware as well as by NSF grant DGE-1106400. We
are grateful to our anonymous reviewers and our shepherd, Miguel Castro,
for thoughtful feedback, which helped improve the quality of this paper.

\section{References}\label{references}

{[}1{]} Akka. https://akka.io/.\\
{[}2{]} Apache Arrow. https://arrow.apache.org/.\\
{[}3{]} Dask Benchmarks.
http://matthewrocklin.com/blog/work/2017/07/03/scaling.\\
{[}4{]} EC2 Instance Pricing.
https://aws.amazon.com/ec2/pricing/on-demand/.\\
{[}5{]} OpenAI Baselines: high-quality implementations of reinforcement
learning algorithms. https://github.com/openai/baselines.\\
{[}6{]} TensorFlow Serving. https://www.tensorflow.org/serving/.\\
{[}7{]} ABADI, M., BARHAM, P., CHEN, J., CHEN, Z., DAVIS, A., DEAN, J.,
DEVIN, M., GHEMAWAT, S., IRVING, G., ISARD, M., ET AL. TensorFlow: A
system for large-scale machine learning. In Proceedings of the 12th
USENIX Symposium on Operating Systems Design and Implementation (OSDI).
Savannah, Georgia, USA (2016).\\
{[}8{]} AGARWAL, A., BIRD, S., COZOWICZ, M., HOANG, L., LANGFORD, J.,
LEE, S., LI, J., MELAMED, D., OSHRI, G., RIBAS, O., SEN, S., AND
SLIVKINS, A. A multiworld testing decision service. arXiv preprint
arXiv:1606.03966 (2016).\\
{[}9{]} ALVARO, P., CONDIE, T., CONWAY, N., ELMELEEGY, K., HELLERSTEIN,
J. M., AND SEARS, R. BOOM Analytics: exploring data-centric, declarative
programming for the cloud. In Proceedings of the 5th European conference
on Computer systems (2010), ACM, pp.~223--236.\\
{[}10{]} ARMSTRONG, J., VIRDING, R., WIKSTRÖM, C., AND WILLIAMS, M.
Concurrent programming in ERLANG.\\
{[}11{]} BEATTIE, C., LEIBO, J. Z., TEPLYASHIN, D., WARD, T.,
WAINWRIGHT, M., KÜTTLER, H., LEFRANCQ, A., GREEN, S., VALDÉS, V., SADIK,
A., ET AL. DeepMind Lab. arXiv preprint arXiv:1612.03801 (2016).\\
{[}12{]} BLUMOFE, R. D., AND LEISERSON, C. E. Scheduling multithreaded
computations by work stealing. J. ACM 46, 5 (Sept.~1999), 720--748.\\
{[}13{]} BROCKMAN, G., CHEUNG, V., PETTERSSON, L., SCHNEIDER, J.,
SCHULMAN, J., TANG, J., AND ZAREMBA, W. OpenAI gym. arXiv preprint
arXiv:1606.01540 (2016).\\
{[}14{]} BYKOV, S., GELLER, A., KLIOT, G., LARUS, J. R., PANDYA, R., AND
THELIN, J. Orleans: Cloud computing for everyone. In Proceedings of the
2nd ACM Symposium on Cloud Computing (2011), ACM, p.~16.\\
{[}15{]} CARBONE, P., EWEN, S., FÓRA, G., HARIDI, S., RICHTER, S., AND
TZOUMAS, K. State management in Apache Flink: Consistent stateful
distributed stream processing. Proc. VLDB Endow. 10, 12 (Aug.~2017),
1718--1729.\\
{[}16{]} CASADO, M., FREEDMAN, M. J., PETTIT, J., LUO, J., MCKEOWN, N.,
AND SHENKER, S. Ethane: Taking control of the enterprise. SIGCOMM
Comput. Commun. Rev.~37, 4 (Aug.~2007), 1--12.\\
{[}17{]} CHAROUSSET, D., SCHMIDT, T. C., HIESGEN, R., AND WÄHLISCH, M.
Native actors: A scalable software platform for distributed,
heterogeneous environments. In Proceedings of the 2013 workshop on
Programming based on actors, agents, and decentralized control (2013),
ACM, pp.~87--96.

{[}18{]} CHEN, T., LI, M., LI, Y., LIN, M., WANG, N., WANG, M., XIAO,
T., XU, B., ZHANG, C., AND ZHANG, Z. MXNet: A flexible and efficient
machine learning library for heterogeneous distributed systems. In NIPS
Workshop on Machine Learning Systems (LearningSys'16) (2016).\\
{[}19{]} CRANKSHAW, D., WANG, X., ZHOU, G., FRANKLIN, M. J., GONZALEZ,
J. E., AND STOICA, I. Clipper: A low-latency online prediction serving
system. In 14th USENIX Symposium on Networked Systems Design and
Implementation (NSDI 17) (Boston, MA, 2017), USENIX Association,
pp.~613--627.\\
{[}20{]} DEAN, J., AND GHEMAWAT, S. MapReduce: Simplified data
processing on large clusters. Commun. ACM 51, 1 (Jan.~2008), 107--113.\\
{[}21{]} DENNIS, J. B., AND MISUNAS, D. P. A preliminary architecture
for a basic data-flow processor. In Proceedings of the 2Nd Annual
Symposium on Computer Architecture (New York, NY, USA, 1975), ISCA `75,
ACM, pp.~126--132.\\
{[}22{]} GABRIEL, E., FAGG, G. E., BOSILCA, G., ANGSKUN, T., DONGARRA,
J. J., SQUYRES, J. M., SAHAY, V., KAMBADUR, P., BARRETT, B., LUMSDAINE,
A., CASTAIN, R. H., DANIEL, D. J., GRAHAM, R. L., AND WOODALL, T. S.
Open MPI: Goals, concept, and design of a next generation MPI
implementation. In Proceedings, 11th European PVM/MPI Users' Group
Meeting (Budapest, Hungary, September 2004), pp.~97--104.\\
{[}23{]} GHEMAWAT, S., GOBIOFF, H., AND LEUNG, S.-T. The Google file
system. 29--43.\\
{[}24{]} GONZALEZ, J. E., XIN, R. S., DAVE, A., CRANKSHAW, D., FRANKLIN,
M. J., AND STOICA, I. GraphX: Graph processing in a distributed dataflow
framework. In Proceedings of the 11th USENIX Conference on Operating
Systems Design and Implementation (Berkeley, CA, USA, 2014), OSDI'14,
USENIX Association, pp.~599--613.\\
{[}25{]} GU*, S., HOLLY*, E., LILLICRAP, T., AND LEVINE, S. Deep
reinforcement learning for robotic manipulation with asynchronous
off-policy updates. In IEEE International Conference on Robotics and
Automation (ICRA 2017) (2017).\\
{[}26{]} HINDMAN, B., KONWINSKI, A., ZAHARIA, M., GHODSI, A., JOSEPH, A.
D., KATZ, R., SHENKER, S., AND STOICA, I. Mesos: A platform for
fine-grained resource sharing in the data center. In Proceedings of the
8th USENIX Conference on Networked Systems Design and Implementation
(Berkeley, CA, USA, 2011), NSDI'11, USENIX Association, pp.~295--308.\\
{[}27{]} HORGAN, D., QUAN, J., BUDDEN, D., BARTH-MARON, G., HESSEL, M.,
VAN HASSELT, H., AND SILVER, D. Distributed prioritized experience
replay. International Conference on Learning Representations (2018).\\
{[}28{]} ISARD, M., BUDIU, M., YU, Y., BIRRELL, A., AND FETTERLY, D.
Dryad: Distributed data-parallel programs from sequential building
blocks. In Proceedings of the 2nd ACM SIGOPS/EuroSys European Conference
on Computer Systems 2007 (New York, NY, USA, 2007), EuroSys '07, ACM,
pp.~59--72.\\
{[}29{]} JIA, Y., SHELHAMER, E., DONAHUE, J., KARAYEV, S., LONG, J.,
GIRSHICK, R., GUADARRAMA, S., AND DARRELL, T. Caffe: Convolutional
architecture for fast feature embedding. arXiv preprint arXiv:1408.5093
(2014).\\
{[}30{]} JORDAN, M. I., AND MITCHELL, T. M. Machine learning: Trends,
perspectives, and prospects. Science 349, 6245 (2015), 255--260.

{[}31{]} LEIBIUSKY, J., EISBRUCH, G., AND SIMONASSI, D. Getting Started
with Storm. O'Reilly Media, Inc., 2012.\\
{[}32{]} LI, M., ANDERSEN, D. G., PARK, J. W., SMOLA, A. J., AHMED, A.,
JOSIFOVSKI, V., LONG, J., SHEKITA, E. J., AND SU, B.-Y. Scaling
distributed machine learning with the parameter server. In Proceedings
of the 11th USENIX Conference on Operating Systems Design and
Implementation (Berkeley, CA, USA, 2014), OSDI'14, pp.~583--598.\\
{[}33{]} LOOKS, M., HERRESHOFF, M., HUTCHINS, D., AND NORVIG, P. Deep
learning with dynamic computation graphs. arXiv preprint
arXiv:1702.02181 (2017).\\
{[}34{]} LOW, Y., GONZALEZ, J., KYROLA, A., BICKSON, D., GUESTRIN, C.,
AND HELLERSTEIN, J. GraphLab: A new framework for parallel machine
learning. In Proceedings of the Twenty-Sixth Conference on Uncertainty
in Artificial Intelligence (Arlington, Virginia, United States, 2010),
UAI'10, pp.~340--349.\\
{[}35{]} MALEWICZ, G., AUSTERN, M. H., BIK, A. J., DEHNERT, J. C., HORN,
I., LEISER, N., AND CZAJKOWSKI, G. Pregel: A system for large-scale
graph processing. In Proceedings of the 2010 ACM SIGMOD International
Conference on Management of Data (New York, NY, USA, 2010), SIGMOD '10,
ACM, pp.~135--146.\\
{[}36{]} MNIH, V., BADIA, A. P., MIRZA, M., GRAVES, A., LILLICRAP, T.
P., HARLEY, T., SILVER, D., AND KAVUKCUOGLU, K. Asynchronous methods for
deep reinforcement learning. In International Conference on Machine
Learning (2016).\\
{[}37{]} MNIH, V., KAVUKCUOGLU, K., SILVER, D., RUSU, A. A., VENESS, J.,
BELLEMARE, M. G., GRAVES, A., RIEDMILLER, M., FIDJELAND, A. K.,
OSTROVSKI, G., ET AL. Human-level control through deep reinforcement
learning. Nature 518, 7540 (2015), 529--533.\\
{[}38{]} MURRAY, D. A Distributed Execution Engine Supporting
Data-dependent Control Flow. University of Cambridge, 2012.\\
{[}39{]} MURRAY, D. G., MCSHERRY, F., ISAACS, R., ISARD, M., BARHAM, P.,
AND ABADI, M. Naiad: A timely dataflow system. In Proceedings of the
Twenty-Fourth ACM Symposium on Operating Systems Principles (New York,
NY, USA, 2013), SOSP '13, ACM, pp.~439--455.\\
{[}40{]} MURRAY, D. G., SCHWARZKOPF, M., SMOWTON, C., SMITH, S.,
MADHAVAPEDDY, A., AND HAND, S. CIEL: A universal execution engine for
distributed data-flow computing. In Proceedings of the 8th USENIX
Conference on Networked Systems Design and Implementation (Berkeley, CA,
USA, 2011), NSDI'11, USENIX Association, pp.~113--126.\\
{[}41{]} NAIR, A., SRINIVASAN, P., BLACKWELL, S., ALCICEK, C., FEARON,
R., MARIA, A. D., PANNEERSHELVAM, V., SULEYMAN, M., BEATTIE, C.,
PETERSEN, S., LEGG, S., MNIH, V., KAVUKCUOGLU, K., AND SILVER, D.
Massively parallel methods for deep reinforcement learning, 2015.\\
{[}42{]} NG, A., COATES, A., DIEL, M., GANAPATHI, V., SCHULTE, J., TSE,
B., BERGER, E., AND LIANG, E. Autonomous inverted helicopter flight via
reinforcement learning. Experimental Robotics IX (2006), 363--372.\\
{[}43{]} NISHIHARA, R., MORITZ, P., WANG, S., TUMANOV, A., PAUL, W.,
SCHLEIER-SMITH, J., LIAW, R., NIKNAMI, M., JORDAN, M. I., AND STOICA, I.
Real-time machine learning: The missing pieces. In Workshop on Hot
Topics in Operating Systems (2017).\\
{[}44{]} OPENAI. OpenAI Dota 2 1v1 bot.
https://openai.com/the-international/, 2017.

{[}45{]} OUSTERHOUT, K., WENDELL, P., ZAHARIA, M., AND STOICA, I.
Sparrow: Distributed, low latency scheduling. In Proceedings of the
Twenty-Fourth ACM Symposium on Operating Systems Principles (New York,
NY, USA, 2013), SOSP '13, ACM, pp.~69--84.\\
{[}46{]} PASZKE, A., GROSS, S., CHINTALA, S., CHANAN, G., YANG, E.,
DEVITO, Z., LIN, Z., DESMAISON, A., ANTIGA, L., AND LERER, A. Automatic
differentiation in PyTorch.\\
{[}47{]} QU, H., MASHAYEKHI, O., TEREI, D., AND LEVIS, P. Canary: A
scheduling architecture for high performance cloud computing. arXiv
preprint arXiv:1602.01412 (2016).\\
{[}48{]} ROCKLIN, M. Dask: Parallel computation with blocked algorithms
and task scheduling. In Proceedings of the 14th Python in Science
Conference (2015), K. Huff and J. Bergstra, Eds., pp.~130 - 136.\\
{[}49{]} SALIMANS, T., HO, J., CHEN, X., AND SUTSKEVER, I. Evolution
strategies as a scalable alternative to reinforcement learning. arXiv
preprint arXiv:1703.03864 (2017).\\
{[}50{]} SANFILIPPO, S. Redis: An open source, in-memory data structure
store. https://redis.io/, 2009.\\
{[}51{]} SCHULMAN, J., WOLSKI, F., DHARIWAL, P., RADFORD, A., AND
KLIMOV, O. Proximal policy optimization algorithms. arXiv preprint
arXiv:1707.06347 (2017).\\
{[}52{]} SCHWARZKOPF, M., KONWINSKI, A., ABD-EL-MALEK, M., AND WILKES,
J. Omega: Flexible, scalable schedulers for large compute clusters. In
Proceedings of the 8th ACM European Conference on Computer Systems (New
York, NY, USA, 2013), EuroSys '13, ACM, pp.~351--364.\\
{[}53{]} SERGEEV, A., AND DEL BALSO, M. Horovod: fast and easy
distributed deep learning in tensorflow. arXiv preprint arXiv:1802.05799
(2018).\\
{[}54{]} SILVER, D., HUANG, A., MADDISON, C. J., GUEZ, A., SIFRE, L.,
VAN DEN DRIESSCHE, G., SCHRITTWIESER, J., ANTONOGLOU, I.,
PANNEERSHELVAM, V., LANCTOT, M., ET AL. Mastering the game of Go with
deep neural networks and tree search. Nature 529, 7587 (2016),
484--489.\\
{[}55{]} SILVER, D., LEVER, G., HEESS, N., DEGRIS, T., WIERSTRA, D., AND
RIEDMILLER, M. Deterministic policy gradient algorithms. In ICML
(2014).\\
{[}56{]} SUTTON, R. S., AND BARTO, A. G. Reinforcement Learning: An
Introduction. MIT press Cambridge, 1998.\\
{[}57{]} THAKUR, R., RABENSEIFNER, R., AND GROPP, W. Optimization of
collective communication operations in MPICH. The International Journal
of High Performance Computing Applications 19, 1 (2005), 49--66.\\
{[}58{]} TIAN, Y., GONG, Q., SHANG, W., WU, Y., AND ZITNICK, C. L. ELF:
An extensive, lightweight and flexible research platform for real-time
strategy games. Advances in Neural Information Processing Systems (NIPS)
(2017).\\
{[}59{]} TODOROV, E., EREZ, T., AND TASSA, Y. Mujoco: A physics engine
for model-based control. In Intelligent Robots and Systems (IROS), 2012
IEEE/RSJ International Conference on (2012), IEEE, pp.~5026--5033.

{[}60{]} VAN DEN BERG, J., MILLER, S., DUCKWORTH, D., HU, H., WAN, A.,
FU, X.-Y., GOLDBERG, K., AND ABBEEL, P. Superhuman performance of
surgical tasks by robots using iterative learning from human-guided
demonstrations. In Robotics and Automation (ICRA), 2010 IEEE
International Conference on (2010), IEEE, pp.~2074--2081.\\
{[}61{]} VAN RENESSE, R., AND SCHNEIDER, F. B. Chain replication for
supporting high throughput and availability. In Proceedings of the 6th
Conference on Symposium on Opearting Systems Design \& Implementation -
Volume 6 (Berkeley, CA, USA, 2004), OSDI'04, USENIX Association.\\
{[}62{]} VENKATARAMAN, S., PANDA, A., OUSTERHOUT, K., GHODSI, A.,
ARMBRUST, M., RECHT, B., FRANKLIN, M., AND STOICA, I. Drizzle: Fast and
adaptable stream processing at scale. In Proceedings of the Twenty-Sixth
ACM Symposium on Operating Systems Principles (2017), SOSP '17, ACM.\\
{[}63{]} WHITE, T. Hadoop: The Definitive Guide. O'Reilly Media, Inc.,
2012.\\
{[}64{]} ZAHARIA, M., CHOWDHURY, M., DAS, T., DAVE, A., MA, J.,
MCCAULEY, M., FRANKLIN, M. J., SHENKER, S., AND STOICA, I. Resilient
distributed datasets: A fault-tolerant abstraction for in-memory cluster
computing. In Proceedings of the 9th USENIX conference on Networked
Systems Design and Implementation (2012), USENIX Association,
pp.~2--2.\\
{[}65{]} ZAHARIA, M., XIN, R. S., WENDELL, P., DAS, T., ARMBRUST, M.,
DAVE, A., MENG, X., ROSEN, J., VENKATARAMAN, S., FRANKLIN, M. J.,
GHODSI, A., GONZALEZ, J., SHENKER, S., AND STOICA, I. Apache Spark: A
unified engine for big data processing. Commun. ACM 59, 11 (Oct.~2016),
56--65.

\end{document}
